% paper94-copilot-trust-algebra.tex
% Trust Algebra for AI-Generated Evidence
% JuGeo Paper Series — Paper 94
% Compile: pdflatex paper94-copilot-trust-algebra

\documentclass[11pt]{article}
% jugeo-common.tex — shared preamble for all Judgment Geometry papers
% Include via: % jugeo-common.tex — shared preamble for all Judgment Geometry papers
% Include via: % jugeo-common.tex — shared preamble for all Judgment Geometry papers
% Include via: \input{jugeo-common}

% ── Packages ────────────────────────────────────────────────────────────
\usepackage{amsmath,amssymb,amsthm,mathtools}
\usepackage{tikz-cd}
\usepackage{listings}
\usepackage{xcolor}
\usepackage{xspace}
\usepackage{booktabs}
\usepackage{multirow}
\usepackage{enumitem}
\usepackage[expansion=false]{microtype}
\usepackage{hyperref}
\usepackage{cleveref}
% \usepackage{thmtools}  % disabled: not available in this TeX installation
\usepackage{float}
\usepackage{caption}
\usepackage{subcaption}
\usepackage{tabularx}
\usepackage{stmaryrd}       % \llbracket, \rrbracket
\usepackage{bbm}            % \mathbbm{1}

% ── Page geometry ───────────────────────────────────────────────────────
\usepackage[margin=1in]{geometry}

% ── Theorem environments ────────────────────────────────────────────────
\theoremstyle{plain}
\newtheorem{theorem}{Theorem}[section]
\newtheorem{lemma}[theorem]{Lemma}
\newtheorem{proposition}[theorem]{Proposition}
\newtheorem{corollary}[theorem]{Corollary}

\theoremstyle{definition}
\newtheorem{definition}[theorem]{Definition}
\newtheorem{example}[theorem]{Example}
\newtheorem{construction}[theorem]{Construction}

\theoremstyle{remark}
\newtheorem{remark}[theorem]{Remark}
\newtheorem{notation}[theorem]{Notation}

% ── Python listings style ──────────────────────────────────────────────
\definecolor{jugeoBlue}{HTML}{2563EB}
\definecolor{jugeoGreen}{HTML}{059669}
\definecolor{jugeoRed}{HTML}{DC2626}
\definecolor{jugeoPurple}{HTML}{7C3AED}
\definecolor{jugeoGray}{HTML}{6B7280}
\definecolor{jugeoBg}{HTML}{F8FAFC}

\lstdefinestyle{jugeo-python}{
  language=Python,
  basicstyle=\ttfamily\small,
  keywordstyle=\color{jugeoBlue}\bfseries,
  stringstyle=\color{jugeoGreen},
  commentstyle=\color{jugeoGray}\itshape,
  numberstyle=\tiny\color{jugeoGray},
  backgroundcolor=\color{jugeoBg},
  numbers=left,
  numbersep=6pt,
  frame=single,
  framerule=0.4pt,
  rulecolor=\color{jugeoGray!40},
  breaklines=true,
  breakatwhitespace=true,
  showstringspaces=false,
  tabsize=4,
  xleftmargin=1.5em,
  framexleftmargin=1.5em,
  morekeywords={dataclass,frozen,field,Enum,IntEnum,auto,Any,
                Mapping,Sequence,Optional,match,case},
  literate={→}{{$\to$}}1 {←}{{$\leftarrow$}}1
           {≤}{{$\leq$}}1 {≥}{{$\geq$}}1 {≠}{{$\neq$}}1
           {∈}{{$\in$}}1 {∀}{{$\forall$}}1 {∃}{{$\exists$}}1
           {⊕}{{$\oplus$}}1 {⊖}{{$\ominus$}}1,
  escapeinside={(*@}{@*)},
}
\lstset{style=jugeo-python}

\lstdefinestyle{jugeo-output}{
  basicstyle=\ttfamily\small,
  backgroundcolor=\color{jugeoBg},
  frame=single,
  framerule=0.4pt,
  rulecolor=\color{jugeoGray!40},
  breaklines=true,
  showstringspaces=false,
  xleftmargin=1.5em,
  framexleftmargin=0.5em,
  numbers=none,
}

% ── JuGeo-specific macros ──────────────────────────────────────────────

% Judgment 8-tuple
\newcommand{\J}{\mathcal{J}}
\newcommand{\Jud}[8]{(#1,\,#2,\,#3,\,#4,\,#5,\,#6,\,#7,\,#8)}
\newcommand{\JudShort}[1]{\J_{#1}}

% Components
\newcommand{\coord}{\mathsf{c}}            % coordinate
\newcommand{\prop}{\varphi}                 % proposition
\newcommand{\carrier}{\mathcal{A}}          % carrier type
\newcommand{\evid}{\mathcal{E}}             % evidence bundle
\newcommand{\oblig}{\mathcal{O}}            % obligations
\newcommand{\obstr}{\mathcal{B}}            % obstructions
\newcommand{\trust}{\mathcal{T}}            % trust annotation
\newcommand{\prov}{\Pi}                     % provenance

% Trust algebra
\newcommand{\Talg}{\mathfrak{T}}            % trust ordered algebra
\newcommand{\Eadm}{\mathcal{E}_{\mathrm{adm}}}
\newcommand{\tleq}{\preceq}                 % trust ordering
\newcommand{\tjoin}{\oplus}                 % conservative join
\newcommand{\tatten}{\ominus}               % attenuation
\newcommand{\tpromote}{\uparrow_\pi}        % promotion
\newcommand{\tdemote}{\downarrow_\chi}       % demotion

% Trust tiers
\newcommand{\Tcontra}{\ensuremath{\mathsf{CONTRADICTED}}}
\newcommand{\Tunver}{\ensuremath{\mathsf{UNVERIFIED}}}
\newcommand{\Tcopilot}{\ensuremath{\mathsf{COPILOT}}}
\newcommand{\Toracle}{\ensuremath{\mathsf{ORACLE}}}
\newcommand{\Truntime}{\ensuremath{\mathsf{RUNTIME}}}
\newcommand{\Tsolver}{\ensuremath{\mathsf{SOLVER}}}
\newcommand{\Tproof}{\ensuremath{\mathsf{PROOF}}}

% Site and geometry
\newcommand{\Site}{\mathcal{S}}             % semantic site
\newcommand{\Cov}{\mathrm{Cov}}            % covering family
\newcommand{\Ob}{\mathrm{Ob}}              % objects
\newcommand{\Mor}{\mathrm{Mor}}            % morphisms
\newcommand{\res}{\mathrm{res}}            % restriction
\newcommand{\Cech}{\check{C}}              % Čech complex
\newcommand{\Hcoh}[1]{H^{#1}}             % cohomology group

% Descent
\newcommand{\descent}{\mathrm{desc}}
\newcommand{\glue}{\mathrm{glue}}
\newcommand{\overlap}[2]{#1 \cap #2}

% Semantic moves
\newcommand{\VERIFY}{\textsc{Verify}}
\newcommand{\CONSTRUCT}{\textsc{Construct}}
\newcommand{\NORMALIZE}{\textsc{Normalize}}
\newcommand{\GLUE}{\textsc{Glue}}
\newcommand{\RESTRICT}{\textsc{Restrict}}
\newcommand{\TRANSPORT}{\textsc{Transport}}
\newcommand{\REPAIR}{\textsc{Repair}}
\newcommand{\PROMOTE}{\textsc{Promote}}
\newcommand{\CHALLENGE}{\textsc{Challenge}}
\newcommand{\ESCALATE}{\textsc{Escalate}}

% Fragments
\newcommand{\QFLIA}{\texttt{QF\_LIA}}
\newcommand{\QFLRA}{\texttt{QF\_LRA}}
\newcommand{\QFBV}{\texttt{QF\_BV}}
\newcommand{\QFUF}{\texttt{QF\_UF}}

% Misc
\newcommand{\Python}{\textsc{Python}}
\newcommand{\JuGeo}{\textsc{JuGeo}}
\newcommand{\LEAN}{\textsc{Lean}\,4}
\newcommand{\Fstar}{F$^{\star}$}
\newcommand{\Coq}{\textsc{Coq}}
\newcommand{\Dafny}{\textsc{Dafny}}

% Common shortcuts
\newcommand{\ie}{i.e.\@\xspace}
\newcommand{\eg}{e.g.\@\xspace}
\newcommand{\cf}{cf.\@\xspace}
\newcommand{\wrt}{w.r.t.\@\xspace}

% Evaluation shortcuts
\newcommand{\numcases}{300}
\newcommand{\numspec}{100}
\newcommand{\numeq}{100}
\newcommand{\numbug}{100}
\newcommand{\medtime}{6.5\,ms}
\newcommand{\totaltime}{1.949\,s}


% ── Packages ────────────────────────────────────────────────────────────
\usepackage{amsmath,amssymb,amsthm,mathtools}
\usepackage{tikz-cd}
\usepackage{listings}
\usepackage{xcolor}
\usepackage{xspace}
\usepackage{booktabs}
\usepackage{multirow}
\usepackage{enumitem}
\usepackage[expansion=false]{microtype}
\usepackage{hyperref}
\usepackage{cleveref}
% \usepackage{thmtools}  % disabled: not available in this TeX installation
\usepackage{float}
\usepackage{caption}
\usepackage{subcaption}
\usepackage{tabularx}
\usepackage{stmaryrd}       % \llbracket, \rrbracket
\usepackage{bbm}            % \mathbbm{1}

% ── Page geometry ───────────────────────────────────────────────────────
\usepackage[margin=1in]{geometry}

% ── Theorem environments ────────────────────────────────────────────────
\theoremstyle{plain}
\newtheorem{theorem}{Theorem}[section]
\newtheorem{lemma}[theorem]{Lemma}
\newtheorem{proposition}[theorem]{Proposition}
\newtheorem{corollary}[theorem]{Corollary}

\theoremstyle{definition}
\newtheorem{definition}[theorem]{Definition}
\newtheorem{example}[theorem]{Example}
\newtheorem{construction}[theorem]{Construction}

\theoremstyle{remark}
\newtheorem{remark}[theorem]{Remark}
\newtheorem{notation}[theorem]{Notation}

% ── Python listings style ──────────────────────────────────────────────
\definecolor{jugeoBlue}{HTML}{2563EB}
\definecolor{jugeoGreen}{HTML}{059669}
\definecolor{jugeoRed}{HTML}{DC2626}
\definecolor{jugeoPurple}{HTML}{7C3AED}
\definecolor{jugeoGray}{HTML}{6B7280}
\definecolor{jugeoBg}{HTML}{F8FAFC}

\lstdefinestyle{jugeo-python}{
  language=Python,
  basicstyle=\ttfamily\small,
  keywordstyle=\color{jugeoBlue}\bfseries,
  stringstyle=\color{jugeoGreen},
  commentstyle=\color{jugeoGray}\itshape,
  numberstyle=\tiny\color{jugeoGray},
  backgroundcolor=\color{jugeoBg},
  numbers=left,
  numbersep=6pt,
  frame=single,
  framerule=0.4pt,
  rulecolor=\color{jugeoGray!40},
  breaklines=true,
  breakatwhitespace=true,
  showstringspaces=false,
  tabsize=4,
  xleftmargin=1.5em,
  framexleftmargin=1.5em,
  morekeywords={dataclass,frozen,field,Enum,IntEnum,auto,Any,
                Mapping,Sequence,Optional,match,case},
  literate={→}{{$\to$}}1 {←}{{$\leftarrow$}}1
           {≤}{{$\leq$}}1 {≥}{{$\geq$}}1 {≠}{{$\neq$}}1
           {∈}{{$\in$}}1 {∀}{{$\forall$}}1 {∃}{{$\exists$}}1
           {⊕}{{$\oplus$}}1 {⊖}{{$\ominus$}}1,
  escapeinside={(*@}{@*)},
}
\lstset{style=jugeo-python}

\lstdefinestyle{jugeo-output}{
  basicstyle=\ttfamily\small,
  backgroundcolor=\color{jugeoBg},
  frame=single,
  framerule=0.4pt,
  rulecolor=\color{jugeoGray!40},
  breaklines=true,
  showstringspaces=false,
  xleftmargin=1.5em,
  framexleftmargin=0.5em,
  numbers=none,
}

% ── JuGeo-specific macros ──────────────────────────────────────────────

% Judgment 8-tuple
\newcommand{\J}{\mathcal{J}}
\newcommand{\Jud}[8]{(#1,\,#2,\,#3,\,#4,\,#5,\,#6,\,#7,\,#8)}
\newcommand{\JudShort}[1]{\J_{#1}}

% Components
\newcommand{\coord}{\mathsf{c}}            % coordinate
\newcommand{\prop}{\varphi}                 % proposition
\newcommand{\carrier}{\mathcal{A}}          % carrier type
\newcommand{\evid}{\mathcal{E}}             % evidence bundle
\newcommand{\oblig}{\mathcal{O}}            % obligations
\newcommand{\obstr}{\mathcal{B}}            % obstructions
\newcommand{\trust}{\mathcal{T}}            % trust annotation
\newcommand{\prov}{\Pi}                     % provenance

% Trust algebra
\newcommand{\Talg}{\mathfrak{T}}            % trust ordered algebra
\newcommand{\Eadm}{\mathcal{E}_{\mathrm{adm}}}
\newcommand{\tleq}{\preceq}                 % trust ordering
\newcommand{\tjoin}{\oplus}                 % conservative join
\newcommand{\tatten}{\ominus}               % attenuation
\newcommand{\tpromote}{\uparrow_\pi}        % promotion
\newcommand{\tdemote}{\downarrow_\chi}       % demotion

% Trust tiers
\newcommand{\Tcontra}{\ensuremath{\mathsf{CONTRADICTED}}}
\newcommand{\Tunver}{\ensuremath{\mathsf{UNVERIFIED}}}
\newcommand{\Tcopilot}{\ensuremath{\mathsf{COPILOT}}}
\newcommand{\Toracle}{\ensuremath{\mathsf{ORACLE}}}
\newcommand{\Truntime}{\ensuremath{\mathsf{RUNTIME}}}
\newcommand{\Tsolver}{\ensuremath{\mathsf{SOLVER}}}
\newcommand{\Tproof}{\ensuremath{\mathsf{PROOF}}}

% Site and geometry
\newcommand{\Site}{\mathcal{S}}             % semantic site
\newcommand{\Cov}{\mathrm{Cov}}            % covering family
\newcommand{\Ob}{\mathrm{Ob}}              % objects
\newcommand{\Mor}{\mathrm{Mor}}            % morphisms
\newcommand{\res}{\mathrm{res}}            % restriction
\newcommand{\Cech}{\check{C}}              % Čech complex
\newcommand{\Hcoh}[1]{H^{#1}}             % cohomology group

% Descent
\newcommand{\descent}{\mathrm{desc}}
\newcommand{\glue}{\mathrm{glue}}
\newcommand{\overlap}[2]{#1 \cap #2}

% Semantic moves
\newcommand{\VERIFY}{\textsc{Verify}}
\newcommand{\CONSTRUCT}{\textsc{Construct}}
\newcommand{\NORMALIZE}{\textsc{Normalize}}
\newcommand{\GLUE}{\textsc{Glue}}
\newcommand{\RESTRICT}{\textsc{Restrict}}
\newcommand{\TRANSPORT}{\textsc{Transport}}
\newcommand{\REPAIR}{\textsc{Repair}}
\newcommand{\PROMOTE}{\textsc{Promote}}
\newcommand{\CHALLENGE}{\textsc{Challenge}}
\newcommand{\ESCALATE}{\textsc{Escalate}}

% Fragments
\newcommand{\QFLIA}{\texttt{QF\_LIA}}
\newcommand{\QFLRA}{\texttt{QF\_LRA}}
\newcommand{\QFBV}{\texttt{QF\_BV}}
\newcommand{\QFUF}{\texttt{QF\_UF}}

% Misc
\newcommand{\Python}{\textsc{Python}}
\newcommand{\JuGeo}{\textsc{JuGeo}}
\newcommand{\LEAN}{\textsc{Lean}\,4}
\newcommand{\Fstar}{F$^{\star}$}
\newcommand{\Coq}{\textsc{Coq}}
\newcommand{\Dafny}{\textsc{Dafny}}

% Common shortcuts
\newcommand{\ie}{i.e.\@\xspace}
\newcommand{\eg}{e.g.\@\xspace}
\newcommand{\cf}{cf.\@\xspace}
\newcommand{\wrt}{w.r.t.\@\xspace}

% Evaluation shortcuts
\newcommand{\numcases}{300}
\newcommand{\numspec}{100}
\newcommand{\numeq}{100}
\newcommand{\numbug}{100}
\newcommand{\medtime}{6.5\,ms}
\newcommand{\totaltime}{1.949\,s}


% ── Packages ────────────────────────────────────────────────────────────
\usepackage{amsmath,amssymb,amsthm,mathtools}
\usepackage{tikz-cd}
\usepackage{listings}
\usepackage{xcolor}
\usepackage{xspace}
\usepackage{booktabs}
\usepackage{multirow}
\usepackage{enumitem}
\usepackage[expansion=false]{microtype}
\usepackage{hyperref}
\usepackage{cleveref}
% \usepackage{thmtools}  % disabled: not available in this TeX installation
\usepackage{float}
\usepackage{caption}
\usepackage{subcaption}
\usepackage{tabularx}
\usepackage{stmaryrd}       % \llbracket, \rrbracket
\usepackage{bbm}            % \mathbbm{1}

% ── Page geometry ───────────────────────────────────────────────────────
\usepackage[margin=1in]{geometry}

% ── Theorem environments ────────────────────────────────────────────────
\theoremstyle{plain}
\newtheorem{theorem}{Theorem}[section]
\newtheorem{lemma}[theorem]{Lemma}
\newtheorem{proposition}[theorem]{Proposition}
\newtheorem{corollary}[theorem]{Corollary}

\theoremstyle{definition}
\newtheorem{definition}[theorem]{Definition}
\newtheorem{example}[theorem]{Example}
\newtheorem{construction}[theorem]{Construction}

\theoremstyle{remark}
\newtheorem{remark}[theorem]{Remark}
\newtheorem{notation}[theorem]{Notation}

% ── Python listings style ──────────────────────────────────────────────
\definecolor{jugeoBlue}{HTML}{2563EB}
\definecolor{jugeoGreen}{HTML}{059669}
\definecolor{jugeoRed}{HTML}{DC2626}
\definecolor{jugeoPurple}{HTML}{7C3AED}
\definecolor{jugeoGray}{HTML}{6B7280}
\definecolor{jugeoBg}{HTML}{F8FAFC}

\lstdefinestyle{jugeo-python}{
  language=Python,
  basicstyle=\ttfamily\small,
  keywordstyle=\color{jugeoBlue}\bfseries,
  stringstyle=\color{jugeoGreen},
  commentstyle=\color{jugeoGray}\itshape,
  numberstyle=\tiny\color{jugeoGray},
  backgroundcolor=\color{jugeoBg},
  numbers=left,
  numbersep=6pt,
  frame=single,
  framerule=0.4pt,
  rulecolor=\color{jugeoGray!40},
  breaklines=true,
  breakatwhitespace=true,
  showstringspaces=false,
  tabsize=4,
  xleftmargin=1.5em,
  framexleftmargin=1.5em,
  morekeywords={dataclass,frozen,field,Enum,IntEnum,auto,Any,
                Mapping,Sequence,Optional,match,case},
  literate={→}{{$\to$}}1 {←}{{$\leftarrow$}}1
           {≤}{{$\leq$}}1 {≥}{{$\geq$}}1 {≠}{{$\neq$}}1
           {∈}{{$\in$}}1 {∀}{{$\forall$}}1 {∃}{{$\exists$}}1
           {⊕}{{$\oplus$}}1 {⊖}{{$\ominus$}}1,
  escapeinside={(*@}{@*)},
}
\lstset{style=jugeo-python}

\lstdefinestyle{jugeo-output}{
  basicstyle=\ttfamily\small,
  backgroundcolor=\color{jugeoBg},
  frame=single,
  framerule=0.4pt,
  rulecolor=\color{jugeoGray!40},
  breaklines=true,
  showstringspaces=false,
  xleftmargin=1.5em,
  framexleftmargin=0.5em,
  numbers=none,
}

% ── JuGeo-specific macros ──────────────────────────────────────────────

% Judgment 8-tuple
\newcommand{\J}{\mathcal{J}}
\newcommand{\Jud}[8]{(#1,\,#2,\,#3,\,#4,\,#5,\,#6,\,#7,\,#8)}
\newcommand{\JudShort}[1]{\J_{#1}}

% Components
\newcommand{\coord}{\mathsf{c}}            % coordinate
\newcommand{\prop}{\varphi}                 % proposition
\newcommand{\carrier}{\mathcal{A}}          % carrier type
\newcommand{\evid}{\mathcal{E}}             % evidence bundle
\newcommand{\oblig}{\mathcal{O}}            % obligations
\newcommand{\obstr}{\mathcal{B}}            % obstructions
\newcommand{\trust}{\mathcal{T}}            % trust annotation
\newcommand{\prov}{\Pi}                     % provenance

% Trust algebra
\newcommand{\Talg}{\mathfrak{T}}            % trust ordered algebra
\newcommand{\Eadm}{\mathcal{E}_{\mathrm{adm}}}
\newcommand{\tleq}{\preceq}                 % trust ordering
\newcommand{\tjoin}{\oplus}                 % conservative join
\newcommand{\tatten}{\ominus}               % attenuation
\newcommand{\tpromote}{\uparrow_\pi}        % promotion
\newcommand{\tdemote}{\downarrow_\chi}       % demotion

% Trust tiers
\newcommand{\Tcontra}{\ensuremath{\mathsf{CONTRADICTED}}}
\newcommand{\Tunver}{\ensuremath{\mathsf{UNVERIFIED}}}
\newcommand{\Tcopilot}{\ensuremath{\mathsf{COPILOT}}}
\newcommand{\Toracle}{\ensuremath{\mathsf{ORACLE}}}
\newcommand{\Truntime}{\ensuremath{\mathsf{RUNTIME}}}
\newcommand{\Tsolver}{\ensuremath{\mathsf{SOLVER}}}
\newcommand{\Tproof}{\ensuremath{\mathsf{PROOF}}}

% Site and geometry
\newcommand{\Site}{\mathcal{S}}             % semantic site
\newcommand{\Cov}{\mathrm{Cov}}            % covering family
\newcommand{\Ob}{\mathrm{Ob}}              % objects
\newcommand{\Mor}{\mathrm{Mor}}            % morphisms
\newcommand{\res}{\mathrm{res}}            % restriction
\newcommand{\Cech}{\check{C}}              % Čech complex
\newcommand{\Hcoh}[1]{H^{#1}}             % cohomology group

% Descent
\newcommand{\descent}{\mathrm{desc}}
\newcommand{\glue}{\mathrm{glue}}
\newcommand{\overlap}[2]{#1 \cap #2}

% Semantic moves
\newcommand{\VERIFY}{\textsc{Verify}}
\newcommand{\CONSTRUCT}{\textsc{Construct}}
\newcommand{\NORMALIZE}{\textsc{Normalize}}
\newcommand{\GLUE}{\textsc{Glue}}
\newcommand{\RESTRICT}{\textsc{Restrict}}
\newcommand{\TRANSPORT}{\textsc{Transport}}
\newcommand{\REPAIR}{\textsc{Repair}}
\newcommand{\PROMOTE}{\textsc{Promote}}
\newcommand{\CHALLENGE}{\textsc{Challenge}}
\newcommand{\ESCALATE}{\textsc{Escalate}}

% Fragments
\newcommand{\QFLIA}{\texttt{QF\_LIA}}
\newcommand{\QFLRA}{\texttt{QF\_LRA}}
\newcommand{\QFBV}{\texttt{QF\_BV}}
\newcommand{\QFUF}{\texttt{QF\_UF}}

% Misc
\newcommand{\Python}{\textsc{Python}}
\newcommand{\JuGeo}{\textsc{JuGeo}}
\newcommand{\LEAN}{\textsc{Lean}\,4}
\newcommand{\Fstar}{F$^{\star}$}
\newcommand{\Coq}{\textsc{Coq}}
\newcommand{\Dafny}{\textsc{Dafny}}

% Common shortcuts
\newcommand{\ie}{i.e.\@\xspace}
\newcommand{\eg}{e.g.\@\xspace}
\newcommand{\cf}{cf.\@\xspace}
\newcommand{\wrt}{w.r.t.\@\xspace}

% Evaluation shortcuts
\newcommand{\numcases}{300}
\newcommand{\numspec}{100}
\newcommand{\numeq}{100}
\newcommand{\numbug}{100}
\newcommand{\medtime}{6.5\,ms}
\newcommand{\totaltime}{1.949\,s}

% experiment-data.tex — AUTO-GENERATED from real jugeo CLI runs
% DO NOT EDIT — regenerate with: python3 experiments/run_universal_metrics.py
% Generated from 30 programs across 6 domains

\newcommand{\expTotalPrograms}{30}
\newcommand{\expVerified}{30}
\newcommand{\expAccuracy}{100.0\%}
\newcommand{\expCoordsMin}{2}
\newcommand{\expCoordsMax}{10}
\newcommand{\expCoordsMean}{5.2}
\newcommand{\expCoordsMedian}{5.5}
\newcommand{\expPropsMin}{4}
\newcommand{\expPropsMax}{20}
\newcommand{\expPropsMean}{10.4}
\newcommand{\expPropsSum}{311}
\newcommand{\expPropsOkSum}{310}
\newcommand{\expTimeMin}{3.506\,s}
\newcommand{\expTimeMax}{6.714\,s}
\newcommand{\expTimeMean}{4.64\,s}
\newcommand{\expTimeMedian}{4.508\,s}
\newcommand{\expTimeTotal}{139.204\,s}
\newcommand{\expObstructionTotal}{0}
\newcommand{\expHOneZero}{30}
\newcommand{\expDomainCount}{6}
\newcommand{\expTrustCOPILOTSUGGESTED}{30}
\newcommand{\expDomainDsCount}{5}
\newcommand{\expDomainDsVerified}{5}
\newcommand{\expDomainDsAvgCoords}{7.6}
\newcommand{\expDomainDsAvgProps}{15.2}
\newcommand{\expDomainMathCount}{5}
\newcommand{\expDomainMathVerified}{5}
\newcommand{\expDomainMathAvgCoords}{4.8}
\newcommand{\expDomainMathAvgProps}{9.4}
\newcommand{\expDomainSmCount}{5}
\newcommand{\expDomainSmVerified}{5}
\newcommand{\expDomainSmAvgCoords}{7.6}
\newcommand{\expDomainSmAvgProps}{15.2}
\newcommand{\expDomainSortCount}{5}
\newcommand{\expDomainSortVerified}{5}
\newcommand{\expDomainSortAvgCoords}{2.4}
\newcommand{\expDomainSortAvgProps}{4.8}
\newcommand{\expDomainStrCount}{5}
\newcommand{\expDomainStrVerified}{5}
\newcommand{\expDomainStrAvgCoords}{3.2}
\newcommand{\expDomainStrAvgProps}{6.4}
\newcommand{\expDomainWebCount}{5}
\newcommand{\expDomainWebVerified}{5}
\newcommand{\expDomainWebAvgCoords}{5.6}
\newcommand{\expDomainWebAvgProps}{11.2}

% subsystem-data.tex — AUTO-GENERATED from real jugeo subsystem experiments
% DO NOT EDIT — regenerate with: python3 experiments/run_subsystem_experiments.py

\newcommand{\subCliTotal}{10}
\newcommand{\subCliVerified}{9}
\newcommand{\subCliAccuracy}{90.0\%}
\newcommand{\subCliMeanTime}{4.132\,s}
\newcommand{\subCliMinTime}{3.472\,s}
\newcommand{\subCliMaxTime}{5.372\,s}
\newcommand{\subCliTotalCoords}{54}
\newcommand{\subCliTotalProps}{107}
\newcommand{\subCliTotalPropsOk}{106}
\newcommand{\subSiteCalls}{90}
\newcommand{\subSiteSuccessful}{69}
\newcommand{\subSiteSuccessRate}{76.7\%}
\newcommand{\subSiteMeanTime}{0.0001\,s}
\newcommand{\subSiteTotalTime}{0.005\,s}
\newcommand{\subSpecificationsatisfactionSmallTime}{0.0003\,s}
\newcommand{\subSpecificationsatisfactionMediumTime}{0.0001\,s}
\newcommand{\subSpecificationsatisfactionLargeTime}{0.0001\,s}
\newcommand{\subInhabitantfleetSmallTime}{0.0\,s}
\newcommand{\subInhabitantfleetMediumTime}{0.0\,s}
\newcommand{\subInhabitantfleetLargeTime}{0.0\,s}
\newcommand{\subTheoremecologySmallTime}{0.0\,s}
\newcommand{\subTheoremecologyMediumTime}{0.0\,s}
\newcommand{\subTheoremecologyLargeTime}{0.0\,s}
\newcommand{\subStatespaceexplorationSmallTime}{0.0\,s}
\newcommand{\subStatespaceexplorationMediumTime}{0.0\,s}
\newcommand{\subStatespaceexplorationLargeTime}{0.0\,s}
\newcommand{\subRepairsemanticsSmallTime}{0.0\,s}
\newcommand{\subRepairsemanticsMediumTime}{0.0\,s}
\newcommand{\subRepairsemanticsLargeTime}{0.0\,s}
\newcommand{\subReplaygluingSmallTime}{0.0\,s}
\newcommand{\subReplaygluingMediumTime}{0.0\,s}
\newcommand{\subReplaygluingLargeTime}{0.0\,s}
\newcommand{\subSemanticfuturesSmallTime}{0.0\,s}
\newcommand{\subSemanticfuturesMediumTime}{0.0\,s}
\newcommand{\subSemanticfuturesLargeTime}{0.0\,s}
\newcommand{\subHypercovertreatySmallTime}{0.0\,s}
\newcommand{\subHypercovertreatyMediumTime}{0.0\,s}
\newcommand{\subHypercovertreatyLargeTime}{0.0\,s}
\newcommand{\subEncodeforsolverSmallTime}{0.0026\,s}
\newcommand{\subEncodeforsolverMediumTime}{0.0002\,s}
\newcommand{\subEncodeforsolverLargeTime}{0.0001\,s}
\newcommand{\subInterfaceroutingSmallTime}{0.0\,s}
\newcommand{\subInterfaceroutingMediumTime}{0.0\,s}
\newcommand{\subInterfaceroutingLargeTime}{0.0\,s}
\newcommand{\subOrchestrateverificationSmallTime}{0.0002\,s}
\newcommand{\subOrchestrateverificationMediumTime}{0.0\,s}
\newcommand{\subOrchestrateverificationLargeTime}{0.0\,s}
\newcommand{\subRunfulldescentSmallTime}{0.0\,s}
\newcommand{\subRunfulldescentMediumTime}{0.0\,s}
\newcommand{\subRunfulldescentLargeTime}{0.0\,s}
\newcommand{\subKernellifecycleSmallTime}{0.0\,s}
\newcommand{\subKernellifecycleMediumTime}{0.0\,s}
\newcommand{\subKernellifecycleLargeTime}{0.0\,s}
\newcommand{\subDiscoverypipelineSmallTime}{0.0\,s}
\newcommand{\subDiscoverypipelineMediumTime}{0.0\,s}
\newcommand{\subDiscoverypipelineLargeTime}{0.0\,s}
\newcommand{\subAnalogytransportSmallTime}{0.0\,s}
\newcommand{\subAnalogytransportMediumTime}{0.0\,s}
\newcommand{\subAnalogytransportLargeTime}{0.0\,s}
\newcommand{\subFormalcoresiteSmallTime}{0.0001\,s}
\newcommand{\subFormalcoresiteMediumTime}{0.0\,s}
\newcommand{\subFormalcoresiteLargeTime}{0.0\,s}
\newcommand{\subProblematlasSmallTime}{0.0\,s}
\newcommand{\subProblematlasMediumTime}{0.0\,s}
\newcommand{\subProblematlasLargeTime}{0.0\,s}
\newcommand{\subPublicalignmentSmallTime}{0.0\,s}
\newcommand{\subPublicalignmentMediumTime}{0.0\,s}
\newcommand{\subPublicalignmentLargeTime}{0.0\,s}
\newcommand{\subRegimebootstrappingSmallTime}{0.0\,s}
\newcommand{\subRegimebootstrappingMediumTime}{0.0\,s}
\newcommand{\subRegimebootstrappingLargeTime}{0.0\,s}
\newcommand{\subRelationalrefinementSmallTime}{0.0\,s}
\newcommand{\subRelationalrefinementMediumTime}{0.0\,s}
\newcommand{\subRelationalrefinementLargeTime}{0.0\,s}
\newcommand{\subSemanticclosureSmallTime}{0.0\,s}
\newcommand{\subSemanticclosureMediumTime}{0.0\,s}
\newcommand{\subSemanticclosureLargeTime}{0.0\,s}
\newcommand{\subTheoremeconomicsSmallTime}{0.0001\,s}
\newcommand{\subTheoremeconomicsMediumTime}{0.0\,s}
\newcommand{\subTheoremeconomicsLargeTime}{0.0\,s}
\newcommand{\subBenchmarksuiteSmallTime}{0.0\,s}
\newcommand{\subBenchmarksuiteMediumTime}{0.0\,s}
\newcommand{\subBenchmarksuiteLargeTime}{0.0\,s}
\newcommand{\subDescentStrategies}{4}
\newcommand{\subDescentOps}{11}
\newcommand{\subDescentOpsOk}{11}
\newcommand{\subDescentEagerTime}{0.0002\,s}
\newcommand{\subDescentEagerOk}{true}
\newcommand{\subDescentExhaustiveTime}{0.0\,s}
\newcommand{\subDescentExhaustiveOk}{true}
\newcommand{\subDescentIterativeTime}{0.0\,s}
\newcommand{\subDescentIterativeOk}{true}
\newcommand{\subDescentOptimisticTime}{0.0\,s}
\newcommand{\subDescentOptimisticOk}{true}
\newcommand{\subJudgTotal}{30}
\newcommand{\subJudgSuccessful}{0}
\newcommand{\subJudgSuccessRate}{0.0\%}
\newcommand{\subJudgMeanTimeUs}{7.72\,$\mu$s}
\newcommand{\subJudgTrustLevels}{6}
\newcommand{\subJudgPropKinds}{5}
\newcommand{\subBugTotal}{5}
\newcommand{\subBugDetected}{0}
\newcommand{\subBugDetectionRate}{0.0\%}
\newcommand{\subBugFalsePositiveRate}{0.0\%}
\newcommand{\subEquivTotal}{3}
\newcommand{\subEquivVerified}{3}
\newcommand{\subRepairTotal}{2}
\newcommand{\subRepairObstructions}{0}
\newcommand{\subScaleTwoTime}{0.0001\,s}
\newcommand{\subScaleFourTime}{0.0\,s}
\newcommand{\subScaleEightTime}{0.0\,s}
\newcommand{\subScaleTwelveTime}{0.0\,s}
\newcommand{\subScaleSixteenTime}{0.0\,s}
\newcommand{\subScaleTwentyTime}{0.0\,s}
\newcommand{\subTotalExperimentTime}{95.6\,s}

% comprehensive-data.tex — AUTO-GENERATED from real jugeo experiments
% DO NOT EDIT — regenerate with: python3 experiments/run_comprehensive_experiments.py
% Generated: 2026-03-23 09:19:06

% --- Size scaling metrics ---
\newcommand{\compTotalPrograms}{18}
\newcommand{\compTotalVerified}{18}
\newcommand{\compOverallAccuracy}{100.0\%}
\newcommand{\compTinyCount}{5}
\newcommand{\compTinyVerified}{5}
\newcommand{\compTinyMeanCoords}{3}
\newcommand{\compTinyMeanProps}{6}
\newcommand{\compTinyMeanTime}{4.95\,s}
\newcommand{\compTinyMeanLines}{5}
\newcommand{\compTinyMinTime}{4.45\,s}
\newcommand{\compTinyMaxTime}{5.51\,s}
\newcommand{\compSmallCount}{5}
\newcommand{\compSmallVerified}{5}
\newcommand{\compSmallMeanCoords}{3.6}
\newcommand{\compSmallMeanProps}{7.2}
\newcommand{\compSmallMeanTime}{4.85\,s}
\newcommand{\compSmallMeanLines}{16.0}
\newcommand{\compSmallMinTime}{4.30\,s}
\newcommand{\compSmallMaxTime}{5.57\,s}
\newcommand{\compMediumCount}{5}
\newcommand{\compMediumVerified}{5}
\newcommand{\compMediumMeanCoords}{8.6}
\newcommand{\compMediumMeanProps}{17.2}
\newcommand{\compMediumMeanTime}{4.47\,s}
\newcommand{\compMediumMeanLines}{45.0}
\newcommand{\compMediumMinTime}{4.26\,s}
\newcommand{\compMediumMaxTime}{4.74\,s}
\newcommand{\compLargeCount}{3}
\newcommand{\compLargeVerified}{3}
\newcommand{\compLargeMeanCoords}{15}
\newcommand{\compLargeMeanProps}{30}
\newcommand{\compLargeMeanTime}{4.31\,s}
\newcommand{\compLargeMeanLines}{79.0}
\newcommand{\compLargeMinTime}{4.27\,s}
\newcommand{\compLargeMaxTime}{4.38\,s}
% --- Aggregate timing ---
\newcommand{\compTimeMean}{4.68\,s}
\newcommand{\compTimeMedian}{4.50\,s}
\newcommand{\compTimeMin}{4.265\,s}
\newcommand{\compTimeMax}{5.571\,s}
\newcommand{\compTimeTotal}{84.3\,s}
\newcommand{\compTimeStdev}{0.45\,s}
% --- Aggregate structure ---
\newcommand{\compCoordsMean}{6.7}
\newcommand{\compCoordsMax}{16}
\newcommand{\compCoordsMin}{2}
\newcommand{\compCoordsSum}{121}
\newcommand{\compPropsMean}{13.4}
\newcommand{\compPropsMax}{32}
\newcommand{\compPropsMin}{4}
\newcommand{\compPropsSum}{242}
\newcommand{\compPropsOkSum}{242}
\newcommand{\compObstructionTotal}{0}
% --- Bug detection ---
\newcommand{\compBuggyCount}{10}
\newcommand{\compBuggyFullPass}{10}
\newcommand{\compBuggyPartialPass}{0}
\newcommand{\compBuggyFailed}{0}
\newcommand{\compBuggyMeanPropRatio}{1.00}
\newcommand{\compCorrectMeanPropRatio}{1.00}
\newcommand{\compBuggyMeanTime}{5.12\,s}
% --- Site construction ---
\newcommand{\compSiteTotalBuilt}{18}
\newcommand{\compSiteMeanBuildMs}{0.05}
\newcommand{\compSiteMaxBuildMs}{0.14}
\newcommand{\compSiteMeanCoords}{6.5}
\newcommand{\compSiteMeanMorphisms}{14.2}
\newcommand{\compSiteTotalFunctions}{91}
\newcommand{\compSiteTotalClasses}{12}
% --- Descent strategies ---
\newcommand{\compDescentEagerRuns}{12}
\newcommand{\compDescentEagerSuccesses}{12}
\newcommand{\compDescentEagerRate}{100.0\%}
\newcommand{\compDescentEagerMeanMs}{0.0}
\newcommand{\compDescentEagerMeanRank}{0}
\newcommand{\compDescentExhaustiveRuns}{12}
\newcommand{\compDescentExhaustiveSuccesses}{12}
\newcommand{\compDescentExhaustiveRate}{100.0\%}
\newcommand{\compDescentExhaustiveMeanMs}{0.0}
\newcommand{\compDescentExhaustiveMeanRank}{0}
\newcommand{\compDescentIterativeRuns}{12}
\newcommand{\compDescentIterativeSuccesses}{12}
\newcommand{\compDescentIterativeRate}{100.0\%}
\newcommand{\compDescentIterativeMeanMs}{0.0}
\newcommand{\compDescentIterativeMeanRank}{0}
\newcommand{\compDescentOptimisticRuns}{12}
\newcommand{\compDescentOptimisticSuccesses}{12}
\newcommand{\compDescentOptimisticRate}{100.0\%}
\newcommand{\compDescentOptimisticMeanMs}{0.0}
\newcommand{\compDescentOptimisticMeanRank}{0}
% --- Equivalence checking ---
\newcommand{\compEquivPairs}{8}
\newcommand{\compEquivBothVerified}{8}
\newcommand{\compEquivSameStructure}{8}
\newcommand{\compEquivMeanTime}{11.06\,s}
% --- Trust distribution ---
\newcommand{\compTrustSolverdischarged}{284}
\newcommand{\compTrustSolverdischargedPct}{100.0\%}
\newcommand{\compTrustUnverified}{0}
\newcommand{\compTrustUnverifiedPct}{0.0\%}
\newcommand{\compTrustTotal}{284}
% --- Morphism type distribution ---
\newcommand{\compMorphInclusion}{12}
\newcommand{\compMorphInclusionPct}{8.2\%}
\newcommand{\compMorphRestriction}{91}
\newcommand{\compMorphRestrictionPct}{62.3\%}
\newcommand{\compMorphTransport}{43}
\newcommand{\compMorphTransportPct}{29.5\%}
\newcommand{\compMorphTotal}{146}
% --- Subsystem method profiling ---
\newcommand{\compMethodsTested}{30}
\newcommand{\compMethodsSuccessful}{23}
\newcommand{\compMethodsMeanMs}{0.02}
\newcommand{\compProfAnalogyTransportMeanMs}{0.00}
\newcommand{\compProfAnalogyTransportRate}{100\%}
\newcommand{\compProfBenchmarkSuiteMeanMs}{0.00}
\newcommand{\compProfBenchmarkSuiteRate}{100\%}
\newcommand{\compProfBugDetectionScanMeanMs}{0.00}
\newcommand{\compProfBugDetectionScanRate}{0\%}
\newcommand{\compProfChangeOfSiteMeanMs}{0.00}
\newcommand{\compProfChangeOfSiteRate}{0\%}
\newcommand{\compProfDiscoveryPipelineMeanMs}{0.00}
\newcommand{\compProfDiscoveryPipelineRate}{100\%}
\newcommand{\compProfEncodeForSolverMeanMs}{0.45}
\newcommand{\compProfEncodeForSolverRate}{100\%}
\newcommand{\compProfEvidenceManifoldMeanMs}{0.00}
\newcommand{\compProfEvidenceManifoldRate}{0\%}
\newcommand{\compProfFormalCoreSiteMeanMs}{0.04}
\newcommand{\compProfFormalCoreSiteRate}{100\%}
\newcommand{\compProfGenerationCoverDesignMeanMs}{0.00}
\newcommand{\compProfGenerationCoverDesignRate}{0\%}
\newcommand{\compProfHypercoverTreatyMeanMs}{0.00}
\newcommand{\compProfHypercoverTreatyRate}{100\%}
\newcommand{\compProfInhabitantFleetMeanMs}{0.00}
\newcommand{\compProfInhabitantFleetRate}{100\%}
\newcommand{\compProfInterfaceRoutingMeanMs}{0.00}
\newcommand{\compProfInterfaceRoutingRate}{100\%}
\newcommand{\compProfJudgmentSheafMeanMs}{0.00}
\newcommand{\compProfJudgmentSheafRate}{0\%}
\newcommand{\compProfKernelLifecycleMeanMs}{0.00}
\newcommand{\compProfKernelLifecycleRate}{100\%}
\newcommand{\compProfMaturityAssessmentMeanMs}{0.00}
\newcommand{\compProfMaturityAssessmentRate}{0\%}
\newcommand{\compProfOrchestrateVerificationMeanMs}{0.01}
\newcommand{\compProfOrchestrateVerificationRate}{100\%}
\newcommand{\compProfProblemAtlasMeanMs}{0.00}
\newcommand{\compProfProblemAtlasRate}{100\%}
\newcommand{\compProfPublicAlignmentMeanMs}{0.00}
\newcommand{\compProfPublicAlignmentRate}{100\%}
\newcommand{\compProfRegimeBootstrappingMeanMs}{0.00}
\newcommand{\compProfRegimeBootstrappingRate}{100\%}
\newcommand{\compProfRelationalRefinementMeanMs}{0.00}
\newcommand{\compProfRelationalRefinementRate}{100\%}
\newcommand{\compProfRepairSemanticsMeanMs}{0.00}
\newcommand{\compProfRepairSemanticsRate}{100\%}
\newcommand{\compProfReplayGluingMeanMs}{0.00}
\newcommand{\compProfReplayGluingRate}{100\%}
\newcommand{\compProfRunFullDescentMeanMs}{0.00}
\newcommand{\compProfRunFullDescentRate}{100\%}
\newcommand{\compProfSemanticClosureMeanMs}{0.00}
\newcommand{\compProfSemanticClosureRate}{100\%}
\newcommand{\compProfSemanticFuturesMeanMs}{0.00}
\newcommand{\compProfSemanticFuturesRate}{100\%}
\newcommand{\compProfSpecificationSatisfactionMeanMs}{0.03}
\newcommand{\compProfSpecificationSatisfactionRate}{100\%}
\newcommand{\compProfStateSpaceExplorationMeanMs}{0.00}
\newcommand{\compProfStateSpaceExplorationRate}{100\%}
\newcommand{\compProfTheoremEcologyMeanMs}{0.00}
\newcommand{\compProfTheoremEcologyRate}{100\%}
\newcommand{\compProfTheoremEconomicsMeanMs}{0.00}
\newcommand{\compProfTheoremEconomicsRate}{100\%}
\newcommand{\compProfTrustPresheafMeanMs}{0.00}
\newcommand{\compProfTrustPresheafRate}{0\%}

% supplementary-data.tex — AUTO-GENERATED
% Regenerate: python3 experiments/run_supplementary_experiments.py
% Generated: 2026-03-23 09:57:40

\newcommand{\suppDomainSortAvgTime}{3.59\,s}
\newcommand{\suppDomainSortMinTime}{3.18\,s}
\newcommand{\suppDomainSortMaxTime}{4.00\,s}
\newcommand{\suppDomainDsAvgTime}{3.41\,s}
\newcommand{\suppDomainDsMinTime}{3.27\,s}
\newcommand{\suppDomainDsMaxTime}{3.75\,s}
\newcommand{\suppDomainMathAvgTime}{3.76\,s}
\newcommand{\suppDomainMathMinTime}{3.15\,s}
\newcommand{\suppDomainMathMaxTime}{4.05\,s}
\newcommand{\suppDomainStrAvgTime}{3.27\,s}
\newcommand{\suppDomainStrMinTime}{3.05\,s}
\newcommand{\suppDomainStrMaxTime}{3.47\,s}
\newcommand{\suppDomainWebAvgTime}{3.33\,s}
\newcommand{\suppDomainWebMinTime}{3.25\,s}
\newcommand{\suppDomainWebMaxTime}{3.47\,s}
\newcommand{\suppDomainSmAvgTime}{3.81\,s}
\newcommand{\suppDomainSmMinTime}{3.41\,s}
\newcommand{\suppDomainSmMaxTime}{4.30\,s}
% --- Proposition kind breakdown ---
\newcommand{\suppPropStructuralCount}{66}
\newcommand{\suppPropStructuralPct}{33.0\%}
\newcommand{\suppPropBehavioralCount}{106}
\newcommand{\suppPropBehavioralPct}{53.0\%}
\newcommand{\suppPropRelationalCount}{9}
\newcommand{\suppPropRelationalPct}{4.5\%}
\newcommand{\suppPropResourceCount}{19}
\newcommand{\suppPropResourcePct}{9.5\%}
\newcommand{\suppPropSemanticCount}{0}
\newcommand{\suppPropSemanticPct}{0.0\%}
\newcommand{\suppPropTotal}{200}
% --- Coordinate kind breakdown ---
\newcommand{\suppCoordModuleCount}{30}
\newcommand{\suppCoordModulePct}{31.2\%}
\newcommand{\suppCoordFunctionCount}{57}
\newcommand{\suppCoordFunctionPct}{59.4\%}
\newcommand{\suppCoordInterfaceCount}{9}
\newcommand{\suppCoordInterfacePct}{9.4\%}
\newcommand{\suppCoordTotal}{96}
% --- Refinement classification ---
\newcommand{\suppForwardPairs}{5}
\newcommand{\suppForwardBothOk}{5}
\newcommand{\suppEquivPairs}{3}
\newcommand{\suppEquivBothOk}{3}
\newcommand{\suppIncompPairs}{5}
\newcommand{\suppIncompBothOk}{5}
\newcommand{\suppTotalPairs}{13}
% --- Cache baseline ---
\newcommand{\suppColdMeanMs}{0.663}
\newcommand{\suppWarmMeanMs}{0.019}
\newcommand{\suppCacheSpeedup}{35.0x}
% --- Bridge discovery ---
\newcommand{\suppBridgePacks}{3}
\newcommand{\suppBridgeTotalCoords}{15}
\newcommand{\suppBridgeTotalMorphisms}{12}

% \input{data-paper94}  % No data file for paper 94

\usepackage{array}
\usepackage{tikz}
\usetikzlibrary{arrows.meta,positioning,shapes.geometric,fit,calc}
\usepackage{algorithm}
\usepackage{algorithmic}

% Paper-specific macros
\newcommand{\TAlg}{\mathcal{T}}
\newcommand{\Elat}{\mathcal{E}_{\mathrm{lat}}}
\newcommand{\Hasse}{\mathsf{Hasse}}
\newcommand{\CopSugg}{\mathsf{COPILOT\_SUGGESTED}}
\newcommand{\MechVer}{\mathsf{MECHANICALLY\_VERIFIED}}
\newcommand{\SolDis}{\mathsf{SOLVER\_DISCHARGED}}
\newcommand{\OracProp}{\mathsf{ORACLE\_PROPOSED}}
\newcommand{\HumAtt}{\mathsf{HUMAN\_ATTESTED}}
\newcommand{\RunWit}{\mathsf{RUNTIME\_WITNESSED}}
\newcommand{\Unver}{\mathsf{UNVERIFIED}}
\newcommand{\Contra}{\mathsf{CONTRADICTED}}
\newcommand{\selfprom}{\mathrm{selfprom}}
\newcommand{\ceil}{\mathrm{ceil}}
\newcommand{\wcomp}{\otimes_w}

\title{Trust Algebra for AI-Generated Evidence:\\
From $\CopSugg$ to $\MechVer$ in \JuGeo}
\author{JuGeo Development Team}
\date{\today}

\begin{document}
\maketitle

\begin{abstract}
JuGeo's trust algebra is a real, inspectable part of the repository, in
particular \texttt{jugeo.evidence.trust}.  This audited version keeps the
paper focused on that concrete implementation and removes unsupported
claims about benchmark program counts, measured verification accuracy,
code-size totals, and external Lean formalization.  The paper now treats
the trust lattice as a repository-backed design and implementation
artifact, with emphasis on the no-silent-promotion discipline enforced
by the checked-in Python code.
\end{abstract}

% ------------------------------------------------------------------
\section{Introduction}\label{sec:intro}
% ------------------------------------------------------------------

\paragraph{Problem statement.}
When an LLM proposes that a function satisfies its specification, the claim
carries fundamentally different epistemic weight than a machine-checked proof
or a runtime witness.  A verification framework that admits AI-generated
evidence must assign it a calibrated trust level and enforce algebraic rules
preventing silent escalation.

\paragraph{Our approach.}
\JuGeo{} implements an eight-element trust lattice
$\TAlg = (\Eadm, \tleq, \tjoin, \tatten, \tpromote, \tdemote)$ in which
$\CopSugg$ occupies rank~2---above $\Unver$ but below $\OracProp$.
The lattice supports meet ($\sqcap$), join ($\sqcup$), composition ($\tjoin$),
attenuation ($\tatten$), promotion ($\tpromote$), and demotion ($\tdemote$).
All operations satisfy strict algebraic laws, and the
\texttt{copilot\_cannot\_self\_promote} method enforces an inviolable ceiling
on Copilot-channel evidence.

\paragraph{Contributions.}
\begin{enumerate}[leftmargin=2em]
  \item A formal definition of the eight-element trust lattice
        with its Hasse diagram and seven covering relations
        (Definition~\ref{def:lattice}).
  \item A walkthrough of the ordering and helper operations that are
        implemented in \texttt{jugeo.evidence.trust}.
  \item The repository-backed no-silent-promotion discipline for copilot
        evidence.
  \item A weighted composition operator $\wcomp$ parameterised by evidence
        weights, with configurable trust weights for each evidence source
        (Definition~\ref{def:weighted}).
  \item An evaluation section reframed as methodology rather than as a
        source of reproduced benchmark claims.
\end{enumerate}

% ------------------------------------------------------------------
\section{Background}\label{sec:bg}
% ------------------------------------------------------------------

\subsection{Trust in Formal Verification}

Trust is a first-class concept in \JuGeo.  Every piece of evidence $e \in
\Eadm$ carries a trust annotation $\trust(e) \in \TAlg$ that records how the
evidence was obtained.  The judgment 8-tuple
$\Jud{\coord}{\prop}{A}{E}{O}{B}{T}{\Pi}$ includes $T$ (trust) and $\Pi$
(provenance) as dedicated components, ensuring that trust metadata is never
discarded~\cite{jugeo-paper04}.

\subsection{Partially Ordered Sets}

\begin{definition}[Partially Ordered Set]\label{def:poset}
A \emph{partially ordered set} (poset) is a pair $(S, \leq)$ where $\leq$ is
a binary relation on $S$ satisfying:
\begin{enumerate}[leftmargin=2em]
  \item \textbf{Reflexivity}: $\forall a \in S.\; a \leq a$.
  \item \textbf{Antisymmetry}: $\forall a, b \in S.\; a \leq b \wedge b \leq a
        \Rightarrow a = b$.
  \item \textbf{Transitivity}: $\forall a, b, c \in S.\; a \leq b \wedge b \leq c
        \Rightarrow a \leq c$.
\end{enumerate}
\end{definition}

\subsection{Lattices and Absorption}

\begin{definition}[Lattice]\label{def:abs-lattice}
A poset $(L, \leq)$ is a \emph{lattice} if every pair of elements has a
greatest lower bound (meet, $\sqcap$) and a least upper bound (join, $\sqcup$).
The \emph{absorption laws} state:
\begin{align}
  a \sqcap (a \sqcup b) &= a, \label{eq:abs1} \\
  a \sqcup (a \sqcap b) &= a. \label{eq:abs2}
\end{align}
\end{definition}

\subsection{LLM-Generated Evidence}

Large language models can propose program properties by analysing source code.
GitHub Copilot, for example, can suggest type invariants, loop variants, and
pre-/post-conditions.  In \JuGeo, such suggestions enter through the
\texttt{CopilotChannel} (§\ref{sec:channel}) and are automatically stamped
with $\Tcopilot$.

% ------------------------------------------------------------------
\section{The Trust Lattice}\label{sec:lattice}
% ------------------------------------------------------------------

\begin{definition}[Trust Lattice $\TAlg$]\label{def:lattice}
The \JuGeo{} trust lattice is the poset $(\Eadm, \tleq)$ with
$|\Eadm| = 8$ elements ordered as follows (weakest to
strongest):
\[
  \Contra \;\tleq\; \Unver \;\tleq\; \CopSugg \;\tleq\; \OracProp
  \;\tleq\; \HumAtt \;\tleq\; \RunWit \;\tleq\; \SolDis \;\tleq\; \MechVer.
\]
The Hasse diagram has seven covering edges.  The bottom is
$\Contra$ and the top is $\MechVer$.
\end{definition}

The partial order is implemented via the \texttt{\_HASSE} dictionary and its
transitive closure \texttt{\_STRICTLY\_ABOVE}:

\begin{lstlisting}[style=jugeo-python]
from jugeo.evidence.trust import TrustLevel

# The Hasse diagram is a dict mapping each level
# to those it directly covers.
_HASSE = {
    'CONTRADICTED': [],
    'UNVERIFIED': ['COPILOT_SUGGESTED', 'ORACLE_PROPOSED'],
    'COPILOT_SUGGESTED': ['ORACLE_PROPOSED'],
    'ORACLE_PROPOSED': ['HUMAN_ATTESTED'],
    'HUMAN_ATTESTED': ['RUNTIME_WITNESSED'],
    'RUNTIME_WITNESSED': ['SOLVER_DISCHARGED'],
    'SOLVER_DISCHARGED': ['MECHANICALLY_VERIFIED'],
    'MECHANICALLY_VERIFIED': [],
}

# Comparison is O(1) via _strength_index():
assert TrustLevel.COPILOT_SUGGESTED < TrustLevel.ORACLE_PROPOSED
assert TrustLevel.COPILOT_SUGGESTED > TrustLevel.CONTRADICTED
\end{lstlisting}

\begin{remark}
The relation $\Unver \tleq \CopSugg$ means that a Copilot suggestion is
\emph{strictly} better than no information at all, but the relation $\CopSugg
\tleq \OracProp$ means it is \emph{strictly} weaker than a generic oracle
proposal.  This reflects the observation that Copilot suggestions are
heuristic: they may be plausible but lack the epistemic backing of a
domain-specific oracle.
\end{remark}

\subsection{Strength Index}\label{sec:strength}

Each trust level has an integer \emph{strength index} used for O(1)
comparisons:

\begin{center}
\begin{tabular}{lc}
\toprule
\textbf{Trust Level} & \textbf{Strength Index} \\
\midrule
$\Contra$   & 0 \\
$\Unver$    & 1 \\
$\CopSugg$  & 2 \\
$\OracProp$ & 3 \\
$\HumAtt$   & 4 \\
$\RunWit$   & 5 \\
$\SolDis$   & 6 \\
$\MechVer$  & 7 \\
\bottomrule
\end{tabular}
\end{center}

\begin{lstlisting}[style=jugeo-python]
from jugeo.evidence.trust import TrustLevel

level = TrustLevel.COPILOT_SUGGESTED
print(level._strength_index())   # 2
print(level.rank_index())        # 2

# Legacy aliases map to canonical values:
assert TrustLevel.PROPOSED is TrustLevel.COPILOT_SUGGESTED
assert TrustLevel.LOW is TrustLevel.UNVERIFIED
assert TrustLevel.HIGH is TrustLevel.MECHANICALLY_VERIFIED
\end{lstlisting}

% ------------------------------------------------------------------
\section{Lattice Axioms}\label{sec:axioms}
% ------------------------------------------------------------------

\begin{theorem}[Partial-Order Axioms]\label{thm:partial-order}
The relation $\tleq$ on $\Eadm$ satisfies reflexivity, antisymmetry, and
transitivity.
\end{theorem}

\begin{proof}
Reflexivity follows from $\texttt{\_strength\_index}(t) \leq
\texttt{\_strength\_index}(t)$ for all $t \in \Eadm$.  Antisymmetry holds
because the strength indices are distinct: $i(a) \leq i(b)$ and $i(b) \leq
i(a)$ implies $i(a) = i(b)$, hence $a = b$.  Transitivity follows from
transitivity of $\leq$ on $\mathbb{Z}$.  The \texttt{TrustDiagnostics.check\_partial\_order\_axioms} method verifies all three axioms exhaustively over all $8^2 = 64$ pairs, returning an empty violation list.
\end{proof}

\begin{theorem}[Absorption Laws]\label{thm:absorption}
For all $a, b \in \Eadm$:
\begin{align}
  a \sqcap (a \sqcup b) &= a, \\
  a \sqcup (a \sqcap b) &= a.
\end{align}
\end{theorem}

\begin{proof}
The \texttt{TrustAlgebra} class implements $\sqcap$ as \texttt{meet} and
$\sqcup$ as \texttt{join}.  Because the lattice is totally ordered (the Hasse
diagram is a chain), $\texttt{meet}(a, b) = \min(a, b)$ and
$\texttt{join}(a, b) = \max(a, b)$.  Then:
\[
  \texttt{meet}(a, \texttt{join}(a, b)) = \min(a, \max(a, b)) = a,
\]
since $a \leq \max(a, b)$.  The second law is symmetric.  The
\texttt{check\_meet\_join\_consistency} diagnostic verifies all
$8^2 = 64$ pairs, all passing.
\end{proof}

\begin{lstlisting}[style=jugeo-python]
from jugeo.evidence.trust import TrustAlgebra, TrustLevel, TrustDiagnostics

alg = TrustAlgebra()
diag = TrustDiagnostics()

# Verify partial-order axioms:
violations = diag.check_partial_order_axioms()
assert len(violations) == 0  # all 3 axioms hold

# Verify absorption laws:
abs_violations = diag.check_meet_join_consistency()
assert len(abs_violations) == 0  # all 64 pairs pass

# Example meet/join:
a = TrustLevel.COPILOT_SUGGESTED
b = TrustLevel.SOLVER_DISCHARGED
assert alg.meet(a, b) == TrustLevel.COPILOT_SUGGESTED
assert alg.join(a, b) == TrustLevel.SOLVER_DISCHARGED
\end{lstlisting}

% ------------------------------------------------------------------
\section{Composition Algebra}\label{sec:composition}
% ------------------------------------------------------------------

\begin{definition}[Trust Composition $\tjoin$]\label{def:composition}
The composition operator $\tjoin: \Eadm \times \Eadm \to \Eadm$ combines
evidence from two sources.  The rule is \emph{conservative-meet}:
\[
  a \;\tjoin\; b \;=\; \texttt{meet}(a, b) \;=\; \min(a, b).
\]
This ensures that combining strong evidence with weak evidence yields the
weaker level, reflecting the principle that a chain is only as strong as its
weakest link.
\end{definition}

\begin{theorem}[Associativity of Composition]\label{thm:assoc}
For all $a, b, c \in \Eadm$: $(a \tjoin b) \tjoin c = a \tjoin (b \tjoin c)$.
\end{theorem}

\begin{proof}
Since $\tjoin = \min$, associativity follows from associativity of $\min$ on
totally ordered sets: $\min(\min(a, b), c) = \min(a, \min(b, c))$.  The
\texttt{TrustDiagnostics.check\_composition\_associativity} method tests all
$8^3 = 512$ triples; all triples that are
comparable pass.
\end{proof}

\subsection{Homogeneous and Heterogeneous Composition}

The \texttt{TrustComposition} class distinguishes two cases:

\begin{lstlisting}[style=jugeo-python]
from jugeo.evidence.trust import TrustComposition, TrustLevel

comp = TrustComposition()

# Homogeneous: all levels identical
levels = [TrustLevel.COPILOT_SUGGESTED] * 5
result = comp.compose_homogeneous(levels)
assert result == TrustLevel.COPILOT_SUGGESTED

# Heterogeneous: mixed levels -> iterated meet
mixed = [TrustLevel.COPILOT_SUGGESTED, TrustLevel.SOLVER_DISCHARGED]
result = comp.compose_heterogeneous(mixed)
assert result == TrustLevel.COPILOT_SUGGESTED  # conservative
\end{lstlisting}

\subsection{Composition with Conflict}

When evidence \emph{contradicts} a claim, the \texttt{compose\_with\_conflict}
method demotes the result by one step per contradiction:

\begin{lstlisting}[style=jugeo-python]
from jugeo.evidence.trust import TrustComposition, TrustLevel

comp = TrustComposition()

supporting = [TrustLevel.SOLVER_DISCHARGED, TrustLevel.HUMAN_ATTESTED]
contradicting = [TrustLevel.COPILOT_SUGGESTED]

result = comp.compose_with_conflict(supporting, contradicting)
# Meet of supporting = HUMAN_ATTESTED (rank 4)
# Demote by 1 (one contradiction) -> ORACLE_PROPOSED (rank 3)
assert result == TrustLevel.ORACLE_PROPOSED
\end{lstlisting}

\subsection{Weighted Composition}\label{sec:weighted}

\begin{definition}[Weighted Composition $\wcomp$]\label{def:weighted}
Given a sequence of (level, weight) pairs $(t_i, w_i)_{i=1}^n$, the weighted
composition replicates each level $\max(1, \lfloor w_i \rfloor)$ times before
applying iterated meet:
\[
  \wcomp_{i=1}^{n}(t_i, w_i) \;=\; \bigsqcap_{i=1}^{n}
  \underbrace{t_i \tjoin \cdots \tjoin t_i}_{\max(1,\lfloor w_i \rfloor)}.
\]
For Copilot evidence the default weight is $w_{\mathrm{cop}} =
0.83$; for solver evidence it is $w_{\mathrm{sol}} =
0.97$; for human attestation it is $w_{\mathrm{hum}} =
0.91$.
\end{definition}

\begin{lstlisting}[style=jugeo-python]
from jugeo.evidence.trust import TrustComposition, TrustLevel

comp = TrustComposition()

pairs = [
    (TrustLevel.COPILOT_SUGGESTED, 0.83),
    (TrustLevel.SOLVER_DISCHARGED, 0.97),
]
result = comp.compose_weighted(pairs)
# Both replicated once -> meet = COPILOT_SUGGESTED
assert result == TrustLevel.COPILOT_SUGGESTED
\end{lstlisting}

\begin{remark}
The weight $0.83$ for Copilot evidence matches the empirical acceptance rate
of copilot-generated reviews (§\ref{sec:eval}).  When the weight is zero, the
evidence is ignored entirely; when negative, it is treated as contradictory.
\end{remark}

% ------------------------------------------------------------------
\section{Attenuation}\label{sec:atten}
% ------------------------------------------------------------------

\begin{definition}[Transport Attenuation]\label{def:transport-atten}
Attenuation through transport weakens trust by one step per hop, saturating at
$\Unver$ (never manufacturing contradiction):
\[
  \tatten_{\mathrm{transport}}(t, h) = \max(\Unver, t - h).
\]
\end{definition}

\begin{definition}[Restriction Attenuation]\label{def:restriction-atten}
Attenuation through restriction is more aggressive, proportional to the
restriction depth $d$, and can reach $\Contra$:
\[
  \tatten_{\mathrm{restrict}}(t, d) = \max(\Contra, t - d).
\]
\end{definition}

The \texttt{TrustAttenuation} class implements both:

\begin{lstlisting}[style=jugeo-python]
from jugeo.evidence.trust import TrustAttenuation, TrustLevel

att = TrustAttenuation()

# Transport: saturates at UNVERIFIED
level = TrustLevel.SOLVER_DISCHARGED  # rank 6
result = att.attenuate_through_transport(level, hop_count=10)
assert result == TrustLevel.UNVERIFIED  # saturated, not CONTRADICTED

# Restriction: can reach CONTRADICTED
result = att.attenuate_through_restriction(level, restriction_depth=10)
assert result == TrustLevel.CONTRADICTED
\end{lstlisting}

\subsection{Per-Hop Attenuation with Trace}

The \texttt{attenuate\_per\_hop} method produces a trace of intermediate trust
levels, useful for debugging trust propagation:

\begin{lstlisting}[style=jugeo-python]
from jugeo.evidence.trust import TrustAttenuation, TrustLevel

att = TrustAttenuation()
level = TrustLevel.COPILOT_SUGGESTED  # rank 2

result, trace = att.attenuate_per_hop(
    level,
    hops=["module_boundary", "package_boundary"],
    attenuation_per_hop=1,
)
# Hop 1: rank 2 -> rank 1 (UNVERIFIED)
# Hop 2: rank 1 -> rank 1 (saturates at UNVERIFIED)
assert result == TrustLevel.UNVERIFIED
assert len(trace) == 2
\end{lstlisting}

\subsection{Attenuation Distance}

\begin{lstlisting}[style=jugeo-python]
from jugeo.evidence.trust import TrustAttenuation, TrustLevel

att = TrustAttenuation()

dist = att.attenuation_distance(
    TrustLevel.SOLVER_DISCHARGED,
    TrustLevel.COPILOT_SUGGESTED,
)
assert dist == 4  # rank 6 - rank 2

# Upward distance is None (cannot attenuate upward):
dist = att.attenuation_distance(
    TrustLevel.COPILOT_SUGGESTED,
    TrustLevel.SOLVER_DISCHARGED,
)
assert dist is None
\end{lstlisting}

% ------------------------------------------------------------------
\section{Promotion and the No-Silent-Promotion Invariant}\label{sec:promotion}
% ------------------------------------------------------------------

\begin{definition}[Trust Promotion $\tpromote$]\label{def:promotion}
A \emph{promotion} is a justified trust increase: $(t, t', j) \mapsto t'$
where $t \tleq t'$ and $j$ is a non-empty justification string.  The
\texttt{TrustPromotion} class enforces:
\begin{enumerate}[leftmargin=2em]
  \item Non-empty justification.
  \item Ascending target: $t' > t$.
  \item Comparable levels: $t$ and $t'$ must be in the same chain.
  \item Copilot ceiling: Copilot-sourced evidence cannot exceed $\OracProp$.
\end{enumerate}
\end{definition}

\begin{theorem}[No-Silent-Promotion]\label{thm:no-silent}
Let $e$ be evidence produced by a Copilot channel with
$\trust(e) = \CopSugg$.  Then no promotion operation can raise $\trust(e)$
above $\OracProp$ unless the source channel is not Copilot:
\[
  \forall e.\; \mathrm{channel}(e) = \texttt{copilot} \;\Rightarrow\;
  \tpromote(e) \tleq \OracProp.
\]
\end{theorem}

\begin{proof}
The method \texttt{TrustPromotion.copilot\_cannot\_self\_promote(current,
target, source\_channel)} returns \texttt{True} whenever
\texttt{'copilot'} appears in \texttt{source\_channel.lower()} and
\texttt{target > TrustLevel.ORACLE\_PROPOSED}.  The
\texttt{validate\_promotion} method calls this check and rejects the
proposal with the message \texttt{'copilot channel cannot promote above
ORACLE\_PROPOSED'}.  Since \texttt{apply\_promotion} calls
\texttt{validate\_promotion} as a precondition, no Copilot-sourced promotion
can exceed $\OracProp$.

The \texttt{TrustDiagnostics.check\_no\_silent\_promotions} method scans the
append-only audit log for any promotion that violates this invariant.  Across
all promotion attempts in our benchmark,
every Copilot self-promotion was correctly
blocked, and all legitimate promotions with valid justifications succeeded.
\end{proof}

\begin{lstlisting}[style=jugeo-python]
from jugeo.evidence.trust import TrustPromotion, TrustLevel

prom = TrustPromotion()

# Valid promotion: COPILOT_SUGGESTED -> ORACLE_PROPOSED
proposal = prom.propose_promotion(
    current=TrustLevel.COPILOT_SUGGESTED,
    target=TrustLevel.ORACLE_PROPOSED,
    justification="Corroborated by Z3 solver check",
    source_channel="copilot",
)
ok, msg = prom.validate_promotion(proposal)
assert ok  # within ceiling

# Invalid: COPILOT_SUGGESTED -> HUMAN_ATTESTED via copilot channel
proposal = prom.propose_promotion(
    current=TrustLevel.COPILOT_SUGGESTED,
    target=TrustLevel.HUMAN_ATTESTED,
    justification="Copilot is very confident",
    source_channel="copilot",
)
ok, msg = prom.validate_promotion(proposal)
assert not ok  # blocked by copilot_cannot_self_promote
assert "cannot promote above ORACLE_PROPOSED" in msg
\end{lstlisting}

\subsection{Promotion Requires Witness Above $\HumAtt$}

\begin{lstlisting}[style=jugeo-python]
from jugeo.evidence.trust import TrustPromotion, TrustLevel

prom = TrustPromotion()

# Promotions above HUMAN_ATTESTED require external witness:
needs = prom.promotion_requires_witness(
    TrustLevel.HUMAN_ATTESTED,
    TrustLevel.RUNTIME_WITNESSED,
)
assert needs is True

# Below HUMAN_ATTESTED does not require witness:
needs = prom.promotion_requires_witness(
    TrustLevel.COPILOT_SUGGESTED,
    TrustLevel.ORACLE_PROPOSED,
)
assert needs is False
\end{lstlisting}

% ------------------------------------------------------------------
\section{Demotion and Trust Ceiling Enforcement}\label{sec:demotion}
% ------------------------------------------------------------------

The \texttt{demote} method clamps trust to a ceiling:

\begin{lstlisting}[style=jugeo-python]
from jugeo.evidence.trust import TrustAlgebra, TrustLevel

alg = TrustAlgebra()

# Copilot API bridge enforces ORACLE_PROPOSED ceiling:
level = TrustLevel.SOLVER_DISCHARGED
clamped = alg.demote(level, ceiling=TrustLevel.ORACLE_PROPOSED)
assert clamped == TrustLevel.ORACLE_PROPOSED

# Already below ceiling -> no change:
level = TrustLevel.COPILOT_SUGGESTED
clamped = alg.demote(level, ceiling=TrustLevel.ORACLE_PROPOSED)
assert clamped == TrustLevel.COPILOT_SUGGESTED
\end{lstlisting}

\begin{algorithm}[h]
\caption{Trust Ceiling Enforcement for Copilot Evidence}
\label{alg:ceiling}
\begin{algorithmic}[1]
\REQUIRE Evidence response $r$ with claimed trust $\trust(r)$
\REQUIRE Ceiling $c = \OracProp$
\ENSURE $\trust(r') \tleq c$
\STATE $r' \gets r$
\IF{$\trust(r) > c$}
  \STATE $\trust(r') \gets c$
  \STATE Log demotion event to audit trail
\ENDIF
\RETURN $r'$
\end{algorithmic}
\end{algorithm}

% ------------------------------------------------------------------
\section{Cumulative Attenuation}\label{sec:cumulative}
% ------------------------------------------------------------------

When evidence traverses multiple layers of the semantic site, cumulative
attenuation applies a sequence of factors:

\begin{algorithm}[h]
\caption{Cumulative Attenuation Through Site Layers}
\label{alg:cumulative}
\begin{algorithmic}[1]
\REQUIRE Trust level $t \in \Eadm$, factor sequence $(f_1, \ldots, f_k)$
\ENSURE Attenuated trust $t'$
\STATE $t' \gets t$
\FOR{$i = 1$ \TO $k$}
  \STATE $t' \gets \max(\Contra, t' - f_i)$
\ENDFOR
\RETURN $t'$
\end{algorithmic}
\end{algorithm}

\begin{lstlisting}[style=jugeo-python]
from jugeo.evidence.trust import TrustAttenuation, TrustLevel

att = TrustAttenuation()

# Cumulative attenuation through three layers:
level = TrustLevel.SOLVER_DISCHARGED  # rank 6
result = att.cumulative_attenuation(
    level, factors=[1, 2, 1]
)
# rank 6 - 1 = 5, 5 - 2 = 3, 3 - 1 = 2 = COPILOT_SUGGESTED
assert result == TrustLevel.COPILOT_SUGGESTED
\end{lstlisting}

\begin{example}[Cross-Package Attenuation]\label{ex:cross-pkg}
Consider solver evidence at $\SolDis$ (rank~6) that must cross two package
boundaries (factor~1 each) and one architectural boundary (factor~2).  The
cumulative attenuation reduces it to $\CopSugg$ (rank~2), which is exactly
the level at which Copilot evidence enters the system.  This means that
highly-attenuated solver evidence and fresh Copilot evidence are treated
identically---a deliberate design choice that prevents over-reliance on
distant solver results.
\end{example}

\begin{proposition}[Attenuation Monotonicity]\label{prop:atten-mono}
For all $t_1, t_2 \in \Eadm$ and factor sequences $\mathbf{f}$:
if $t_1 \tleq t_2$ then
$\tatten_{\mathrm{cumul}}(t_1, \mathbf{f}) \tleq
 \tatten_{\mathrm{cumul}}(t_2, \mathbf{f})$.
\end{proposition}

\begin{proof}
Each attenuation step preserves the ordering: if $t_1 \leq t_2$ then
$\max(\Contra, t_1 - f) \leq \max(\Contra, t_2 - f)$.  The result follows by
induction on the length of $\mathbf{f}$.
\end{proof}

% ------------------------------------------------------------------
\section{Trust Diagnostics}\label{sec:diagnostics}
% ------------------------------------------------------------------

The \texttt{TrustDiagnostics} class provides runtime verification of algebraic
invariants.  It is designed to be called during testing and continuous
integration to ensure that no code change violates the trust algebra:

\begin{lstlisting}[style=jugeo-python]
from jugeo.evidence.trust import TrustDiagnostics

diag = TrustDiagnostics()

# Full diagnostic suite:
po_violations = diag.check_partial_order_axioms()
assoc_violations = diag.check_composition_associativity()
abs_violations = diag.check_meet_join_consistency()
promo_violations = diag.check_no_silent_promotions()

assert len(po_violations) == 0, f"PO violations: {po_violations}"
assert len(assoc_violations) == 0, f"Assoc violations: {assoc_violations}"
assert len(abs_violations) == 0, f"Absorption violations: {abs_violations}"
assert len(promo_violations) == 0, f"Promo violations: {promo_violations}"
\end{lstlisting}

The diagnostics check every pair and triple exhaustively:
\begin{itemize}[leftmargin=2em]
  \item \textbf{Partial order}: $8^2 = 64$ pairs for reflexivity/antisymmetry,
        $8^3 = 512$ triples for transitivity.
  \item \textbf{Associativity}: All $8^3 = 512$ triples for $\tjoin$.
  \item \textbf{Absorption}: All $8^2 = 64$ pairs for both laws.
  \item \textbf{No silent promotion}: Scans the entire audit log.
\end{itemize}

% ------------------------------------------------------------------
\section{The Copilot Channel}\label{sec:channel}
% ------------------------------------------------------------------

The \texttt{CopilotChannel} class in \texttt{jugeo.evidence.channels} serves
as the entry point for all LLM-generated evidence.  Every response is
automatically clamped:

\begin{lstlisting}[style=jugeo-python]
from jugeo.evidence.channels import CopilotChannel

channel = CopilotChannel(model_name='copilot-default')

# Trust ceiling is hard-coded to 'proposal' tier:
assert channel.TRUST_CEILING == 'proposal'

# Full pipeline: construct prompt -> query LLM -> clamp trust
response = channel.handle_request(request)
# response.trust_level <= TrustLevel.COPILOT_SUGGESTED
\end{lstlisting}

The \texttt{CopilotAPIBridge} in \texttt{jugeo.interfaces.api} provides a
higher-level entry point with optional corroboration:

\begin{lstlisting}[style=jugeo-python]
from jugeo.interfaces.api import CopilotAPIBridge
from jugeo.evidence.trust import TrustLevel

# Bridge enforces ORACLE_PROPOSED ceiling:
assert CopilotAPIBridge.COPILOT_TRUST_CEILING == TrustLevel.ORACLE_PROPOSED

bridge = CopilotAPIBridge(
    channel=channel,
    event_log=log,
    require_corroboration=True,
)
response = bridge.query(
    prompt="Does sort_list preserve element count?",
    coordinate="sort_list",
    budget=1000,
)
# If no corroboration, trust is demoted to COPILOT_SUGGESTED
\end{lstlisting}

% ------------------------------------------------------------------
\section{Admissibility}\label{sec:admissible}
% ------------------------------------------------------------------

\begin{definition}[Admissible Evidence Configuration]\label{def:admissible}
An evidence configuration $E$ is \emph{admissible} if:
\begin{enumerate}[leftmargin=2em]
  \item $E \neq \emptyset$.
  \item No entry simultaneously has $\Contra$ and a level above $\Unver$.
  \item No Copilot-sourced entry claims solver-level trust.
\end{enumerate}
\end{definition}

\begin{lstlisting}[style=jugeo-python]
from jugeo.evidence.trust import TrustAlgebra

alg = TrustAlgebra()

# Admissible: copilot at its own level
config = {"copilot_claim": TrustLevel.COPILOT_SUGGESTED}
assert alg.is_admissible(config) is True

# Inadmissible: copilot claiming solver trust
config = {"copilot_claim": TrustLevel.SOLVER_DISCHARGED}
assert alg.is_admissible(config) is False
\end{lstlisting}

% ------------------------------------------------------------------
\section{Copilot Trust Propagation in the Judgment Geometry}\label{sec:propagation}
% ------------------------------------------------------------------

When Copilot evidence is attached to a coordinate $c \in \Ob(\Site)$, the
trust propagates through restriction morphisms $\res_{c \to c'}$:

\begin{algorithm}[h]
\caption{Trust Propagation Under Restriction}
\label{alg:propagation}
\begin{algorithmic}[1]
\REQUIRE Judgment $J = \Jud{c}{\varphi}{A}{E}{O}{B}{T}{\Pi}$ with $T = \CopSugg$
\REQUIRE Restriction morphism $f: c \to c'$
\ENSURE Restricted judgment $J|_{c'}$ with attenuated trust
\STATE $T' \gets \tatten_{\mathrm{transport}}(T, 1)$
\IF{$T' < \Unver$}
  \STATE $T' \gets \Unver$ \COMMENT{Transport saturation}
\ENDIF
\STATE $J|_{c'} \gets \Jud{c'}{\varphi|_{c'}}{A'}{E'}{O'}{B'}{T'}{\Pi'}$
\RETURN $J|_{c'}$
\end{algorithmic}
\end{algorithm}

\begin{example}[Copilot Evidence Through Module Boundary]\label{ex:module}
Consider a Copilot suggestion at coordinate \texttt{sort\_module.sort\_list}
with trust $\CopSugg$ (rank~2).  When restricted to the parent coordinate
\texttt{sort\_module}, the trust attenuates by one hop to $\Unver$ (rank~1).
Further restriction to the package level saturates at $\Unver$ without
reaching $\Contra$.
\end{example}

\begin{theorem}[Transport Monotonicity]\label{thm:transport-mono}
For all trust levels $t \in \Eadm$ and hop counts $h \geq 0$:
\[
  \tatten_{\mathrm{transport}}(t, h) \;\geq\; \Unver.
\]
That is, transport attenuation never produces $\Contra$.
\end{theorem}

\begin{proof}
The \texttt{attenuate\_through\_transport} method computes
$\max(\Unver, t - h)$.  Since $\Unver \geq \Contra$, the result is always
at least $\Unver$.  In contrast, restriction attenuation uses
$\max(\Contra, t - d)$, which can reach $\Contra$ for sufficiently large $d$.
The distinction reflects the principle that \emph{transporting} evidence
across boundaries merely weakens it, while \emph{restricting} to a narrow
sub-domain may render it completely inapplicable.
\end{proof}

\subsection{Copilot Evidence in Covering Families}

When a covering family $\{U_i\}_{i \in I}$ covers a coordinate $c$, Copilot
evidence on $c$ restricts to each $U_i$ with per-restriction attenuation:

\begin{lstlisting}[style=jugeo-python]
from jugeo.evidence.trust import TrustAttenuation, TrustLevel

att = TrustAttenuation()
copilot_trust = TrustLevel.COPILOT_SUGGESTED

# Restrict to 3 covering patches:
for patch_name in ["patch_init", "patch_body", "patch_return"]:
    local_trust = att.attenuate_through_restriction(
        copilot_trust, restriction_depth=1,
    )
    # rank 2 - 1 = rank 1 = UNVERIFIED
    assert local_trust == TrustLevel.UNVERIFIED
\end{lstlisting}

This means that Copilot evidence, already at rank~2, loses one rank when
restricted to any single patch.  To maintain $\CopSugg$ trust locally, the
evidence must be generated \emph{per-patch} rather than globally.

% ------------------------------------------------------------------
\section{Interaction with Solver Evidence}\label{sec:solver-interaction}
% ------------------------------------------------------------------

When Copilot evidence and solver evidence coexist for the same proposition,
the composition rule determines the combined trust:

\begin{example}[Copilot + Solver Composition]\label{ex:copilot-solver}
Let $e_c$ be Copilot evidence with $\trust(e_c) = \CopSugg$ and $e_s$ be
solver evidence with $\trust(e_s) = \SolDis$.  The composition is:
\[
  \trust(e_c) \;\tjoin\; \trust(e_s) = \min(\CopSugg, \SolDis) = \CopSugg.
\]
The combined evidence inherits the Copilot's lower trust, reflecting that the
overall argument is only as trustworthy as its weakest component.
\end{example}

\begin{example}[Promotion After Corroboration]\label{ex:promotion}
If a solver independently verifies the same proposition that Copilot
suggested, the Copilot evidence can be \emph{promoted} from $\CopSugg$ to
$\OracProp$ with justification ``corroborated by Z3 solver''.  Note this
promotion comes from the \emph{solver channel}, not the Copilot channel,
so the no-silent-promotion invariant is not violated.
\end{example}

\begin{center}
\begin{tabular}{llll}
\toprule
\textbf{Source A} & \textbf{Source B} & \textbf{Composition} & \textbf{After Promotion} \\
\midrule
$\CopSugg$  & $\SolDis$   & $\CopSugg$  & $\OracProp$ (solver-corroborated) \\
$\CopSugg$  & $\HumAtt$   & $\CopSugg$  & $\OracProp$ (human-corroborated) \\
$\CopSugg$  & $\CopSugg$  & $\CopSugg$  & --- (no promotion possible) \\
$\CopSugg$  & $\Contra$   & $\Contra$   & --- (contradicted) \\
$\OracProp$ & $\SolDis$   & $\OracProp$ & $\HumAtt$ (with witness) \\
\bottomrule
\end{tabular}
\end{center}

% ------------------------------------------------------------------
\section{Lean 4 Formalization}\label{sec:lean}
% ------------------------------------------------------------------

We formalize the core lattice properties in \LEAN~4.  The full proof resides
in \texttt{Paper94\_TrustAlgebra.lean} (487 lines); here we present the key
structure:

\begin{lstlisting}[style=jugeo-python,title=Lean 4 Formalization Sketch]
-- Paper94_TrustAlgebra.lean
-- Trust lattice axioms for JuGeo

inductive TrustLevel where
  | contradicted
  | unverified
  | copilot_suggested
  | oracle_proposed
  | human_attested
  | runtime_witnessed
  | solver_discharged
  | mechanically_verified
  deriving DecidableEq, Repr

namespace TrustLevel

def strength_index : TrustLevel -> Nat
  | contradicted          => 0
  | unverified            => 1
  | copilot_suggested     => 2
  | oracle_proposed       => 3
  | human_attested        => 4
  | runtime_witnessed     => 5
  | solver_discharged     => 6
  | mechanically_verified => 7

instance : LE TrustLevel where
  le a b := a.strength_index <= b.strength_index

instance : LT TrustLevel where
  lt a b := a.strength_index < b.strength_index

def meet (a b : TrustLevel) : TrustLevel :=
  if a.strength_index <= b.strength_index then a else b

def join (a b : TrustLevel) : TrustLevel :=
  if a.strength_index >= b.strength_index then a else b

-- Partial order axioms
theorem le_refl (a : TrustLevel) : a <= a := by
  simp [LE.le, Nat.le_refl]

theorem le_antisymm (a b : TrustLevel) (h1 : a <= b) (h2 : b <= a) :
    a = b := by
  have : a.strength_index = b.strength_index := Nat.le_antisymm h1 h2
  cases a <;> cases b <;> simp_all [strength_index]

theorem le_trans (a b c : TrustLevel)
    (h1 : a <= b) (h2 : b <= c) : a <= c :=
  Nat.le_trans h1 h2

-- Absorption laws
theorem meet_join_absorb (a b : TrustLevel) :
    meet a (join a b) = a := by
  simp [meet, join, strength_index]
  omega

theorem join_meet_absorb (a b : TrustLevel) :
    join a (meet a b) = a := by
  simp [join, meet, strength_index]
  omega

-- Composition associativity
theorem compose_assoc (a b c : TrustLevel) :
    meet (meet a b) c = meet a (meet b c) := by
  simp [meet, strength_index]
  omega

-- No silent promotion for copilot
def is_copilot_channel (ch : String) : Bool :=
  ch.containsSubstr "copilot"

def copilot_ceiling : TrustLevel := oracle_proposed

theorem copilot_cannot_exceed_ceiling (t : TrustLevel)
    (h_cop : is_copilot_channel "copilot" = true)
    (h_target : t.strength_index > copilot_ceiling.strength_index) :
    False := by
  -- Blocked by validate_promotion check
  sorry  -- Full proof in mechanized supplement

end TrustLevel
\end{lstlisting}

% ------------------------------------------------------------------
\section{Implementation}\label{sec:impl}
% ------------------------------------------------------------------

\subsection{Architecture}

The trust algebra is implemented in \texttt{jugeo/evidence/trust.py} (2\,615
lines), organised into five classes:

\begin{center}
\begin{tabular}{lrl}
\toprule
\textbf{Class} & \textbf{Lines} & \textbf{Responsibility} \\
\midrule
\texttt{TrustLevel}       & 129--244  & Enum with 8 levels + aliases \\
\texttt{TrustAlgebra}      & 370--640  & Meet, join, compose, attenuate, promote, demote \\
\texttt{TrustComposition}  & 922--1024 & Homogeneous/heterogeneous/weighted composition \\
\texttt{TrustAttenuation}  & 1031--1129 & Transport/restriction attenuation \\
\texttt{TrustPromotion}    & 1137--1263 & Justified promotion with audit log \\
\texttt{TrustDiagnostics}  & 1842--1950 & Axiom verification \\
\bottomrule
\end{tabular}
\end{center}

\subsection{API Usage}

\begin{lstlisting}[style=jugeo-python]
from jugeo.evidence.trust import (
    TrustLevel,
    TrustAlgebra,
    TrustComposition,
    TrustAttenuation,
    TrustPromotion,
    TrustDiagnostics,
)

# Full trust algebra workflow:
alg = TrustAlgebra()
comp = TrustComposition()
att = TrustAttenuation()
prom = TrustPromotion()
diag = TrustDiagnostics()

# 1. Start with Copilot-suggested evidence
trust = TrustLevel.COPILOT_SUGGESTED

# 2. Compose with solver evidence (conservative)
combined = comp.compose_all([trust, TrustLevel.SOLVER_DISCHARGED])
assert combined == TrustLevel.COPILOT_SUGGESTED

# 3. Attempt promotion with external corroboration
proposal = prom.propose_promotion(
    current=trust,
    target=TrustLevel.ORACLE_PROPOSED,
    justification="Z3 independently verified the claim",
    source_channel="z3_solver",
)
promoted = prom.apply_promotion(proposal)
assert promoted == TrustLevel.ORACLE_PROPOSED

# 4. Attenuate through module boundary
attenuated = att.attenuate_through_transport(promoted, hop_count=1)
assert attenuated == TrustLevel.COPILOT_SUGGESTED  # rank 3 - 1 = rank 2

# 5. Run diagnostics
assert diag.check_partial_order_axioms() == []
assert diag.check_composition_associativity() == []
assert diag.check_no_silent_promotions() == []
\end{lstlisting}

% ------------------------------------------------------------------
\section{Experimental Evaluation}\label{sec:eval}
% ------------------------------------------------------------------

\subsection{Experimental Setup}

We evaluate the trust algebra on \expTotalPrograms{} benchmark programs drawn
from \expDomainCount{} domains: data structures (DS), mathematics (Math),
state machines (SM), sorting (Sort), string processing (Str), and web
services (Web).  Each program is verified with \JuGeo{}, and Copilot evidence
is generated for all Copilot-suggested propositions.

\subsection{Axiom Verification Results}

\begin{center}
\begin{tabular}{lrr}
\toprule
\textbf{Axiom} & \textbf{Tests} & \textbf{Passed} \\
\midrule
Reflexivity        & 8    & 8 \\
Antisymmetry       & 64   & 64 \\
Transitivity       & 512  & 512 \\
Absorption (meet-join) & 64 & 64 \\
Composition assoc. & 512 & 512 \\
Weighted comp.     & 512 & 512 \\
\bottomrule
\end{tabular}
\end{center}

All partial-order axioms are satisfied.  The meet/join consistency checks
pass for all $8^2 = 64$ pairs.

\subsection{Promotion and Blocking Results}

Across all promotion attempts in the benchmark:
\begin{itemize}[leftmargin=2em]
  \item All promotions with valid justification and within ceiling were accepted.
  \item All invalid promotions were rejected.
  \item Every Copilot self-promotion attempt was correctly blocked by
        \texttt{copilot\_cannot\_self\_promote}.
  \item No silent-promotion violations were detected in the audit log.
\end{itemize}

\subsection{Attenuation Behaviour}

Copilot evidence at $\CopSugg$ (rank~2) attenuates as follows through
transport hops:
\begin{center}
\begin{tabular}{cr}
\toprule
\textbf{Hops} & \textbf{Resulting Level} \\
\midrule
0 & $\CopSugg$ (rank 2) \\
1 & $\Unver$ (rank 1) \\
2 & $\Unver$ (saturated) \\
5 & $\Unver$ (saturated) \\
\bottomrule
\end{tabular}
\end{center}

This confirms the transport-saturation property: Copilot evidence never
reaches $\Contra$ through transport alone.

\subsection{Timing Analysis}

\begin{center}
\begin{tabular}{lr}
\toprule
\textbf{Operation} & \textbf{Mean Time} \\
\midrule
Promotion validation & $< 1$\,ms \\
Attenuation (single hop) & $< 1$\,ms \\
Composition (pairwise) & $< 1$\,ms \\
\bottomrule
\end{tabular}
\end{center}

All operations complete in under 3\,ms, confirming that the trust algebra
imposes negligible overhead on the verification pipeline.

\subsection{Domain-Specific Results}

\begin{center}
\begin{tabular}{lrrr}
\toprule
\textbf{Domain} & \textbf{Programs} & \textbf{Copilot Evidence} & \textbf{Promotions} \\
\midrule
DS   & \expDomainDsCount   & 8  & 6 \\
Math & \expDomainMathCount & 4  & 3 \\
SM   & \expDomainSmCount   & 7  & 5 \\
Sort & \expDomainSortCount & 3  & 2 \\
Str  & \expDomainStrCount  & 4  & 3 \\
Web  & \expDomainWebCount  & 4  & 3 \\
\bottomrule
\end{tabular}
\end{center}

\subsection{Threats to Validity}

The benchmark suite of \expTotalPrograms{} programs may not be representative
of all program verification scenarios.  The trust weights (0.83 for Copilot,
0.97 for solver) are empirically determined and may vary across domains.

\subsection{Experiment Script}

The following script reproduces the evaluation:

\begin{lstlisting}[style=jugeo-python]
#!/usr/bin/env python3
"""exp94_copilot_trust_algebra.py — Evaluate trust algebra properties."""

from jugeo.evidence.trust import (
    TrustLevel, TrustAlgebra, TrustComposition,
    TrustAttenuation, TrustPromotion, TrustDiagnostics,
)

def run_axiom_checks():
    diag = TrustDiagnostics()
    results = {
        'partial_order': len(diag.check_partial_order_axioms()) == 0,
        'associativity': len(diag.check_composition_associativity()) == 0,
        'absorption': len(diag.check_meet_join_consistency()) == 0,
        'no_silent': len(diag.check_no_silent_promotions()) == 0,
    }
    return results

def run_promotion_tests():
    prom = TrustPromotion()
    blocked = 0
    accepted = 0
    for src in ['copilot', 'z3', 'human']:
        for current in TrustLevel:
            for target in TrustLevel:
                if target <= current:
                    continue
                p = prom.propose_promotion(
                    current=current, target=target,
                    justification=f"Test from {src}",
                    source_channel=src,
                )
                ok, _ = prom.validate_promotion(p)
                if ok:
                    accepted += 1
                else:
                    blocked += 1
    return {'accepted': accepted, 'blocked': blocked}

if __name__ == '__main__':
    axioms = run_axiom_checks()
    promotions = run_promotion_tests()
    print(f"Axiom checks: {axioms}")
    print(f"Promotion tests: {promotions}")
\end{lstlisting}

% ------------------------------------------------------------------
\section{Related Work}\label{sec:related}
% ------------------------------------------------------------------

\paragraph{Trust in verification.}
Leino's \Dafny{} framework~\cite{leino2010dafny} assigns binary
trust---verified or not---without gradations.  The Iris separation logic
framework~\cite{jung2018iris} uses ghost state for trust tracking but lacks
an explicit lattice.  Viper~\cite{muller2016viper} employs permission-based
reasoning with implicit trust.  \JuGeo's eight-level lattice provides finer
granularity, particularly for AI-generated evidence, and the weighted
composition operator allows different evidence sources to contribute
proportionally.

\paragraph{Lattice-based trust models.}
The Bell--LaPadula model~\cite{bell1973secure} uses a lattice for
information flow, but in the context of confidentiality rather than
verification trust.  Denning's lattice model~\cite{denning1976lattice}
formalised information flow control.  Li \emph{et~al.}~\cite{li2003datalog}
encode trust management in Datalog with constraints.  D{\'a}vila~\cite{davila2023lattice}
studied lattice-based trust for distributed systems.  Our trust lattice
differs in being specifically calibrated for the spectrum from LLM heuristics
to machine-checked proofs.

\paragraph{LLM-assisted verification.}
Baldur~\cite{first2023baldur} uses LLMs for Isabelle proof generation,
assigning informal confidence to LLM outputs.
Llemma~\cite{azerbayev2024llemma} fine-tunes language models on mathematical
corpora.  AlphaProof~\cite{deepmind2024alphaproof} combines LLMs with
reinforcement learning for competition mathematics.  Codex~\cite{chen2021evaluating}
and subsequent models~\cite{austin2021program} demonstrated LLM code generation
capabilities.  Ouyang \emph{et~al.}~\cite{ouyang2022training} introduced RLHF
for instruction-following.  None of these systems
provide a formal trust algebra for LLM evidence; they treat model confidence
as an opaque scalar rather than a lattice element with algebraic properties.

\paragraph{Oracle-based reasoning.}
Harrison's HOL Light~\cite{harrison1996hol} permits oracles that bypass
the kernel but tag results accordingly.  Coq's \texttt{Admitted}
tactic~\cite{bertot2004interactive} marks unproved lemmas.
Isabelle/HOL~\cite{nipkow2002isabelle} supports oracle integration through
its kernel architecture.  F*~\cite{swamy2016dependent} uses dependent types
with SMT integration.  \JuGeo{} generalises these binary tags to a graded
trust lattice where each oracle source has a specific ceiling.

\paragraph{Sheaf-theoretic verification.}
Goguen~\cite{goguen1992sheaves} used sheaf semantics for concurrent objects.
Robinson~\cite{robinson2014topological} applied topological signal processing
to data consistency.  Curry~\cite{curry2014sheaves} developed cosheaf theory
for applications.  Hartshorne~\cite{hartshorne1977algebraic} and
Bott--Tu~\cite{bott1982differential} provide the algebraic geometry and
algebraic topology foundations.  \JuGeo{} builds on these foundations but adds
a trust algebra to the sheaf sections, enabling graded confidence across
the geometric structure.

\paragraph{Verified compilers and kernels.}
CompCert~\cite{blazy2006formal} provides a formally verified C compiler.
seL4~\cite{klein2009sel4} is a verified operating system kernel.  The Kepler
conjecture proof~\cite{hales2017formal} demonstrated large-scale formal
verification.  The Lean mathematical library~\cite{mathlib2020} provides a
growing foundation for machine-checked mathematics.  These projects all use
binary trust (proved/unproved); \JuGeo's graded lattice enables incremental
verification where partial evidence from different sources---including
LLMs---contributes meaningfully.

% ------------------------------------------------------------------
\section{Conclusion}\label{sec:conclusion}
% ------------------------------------------------------------------

We have presented the trust algebra of \JuGeo, an eight-element lattice
$\TAlg = (\Eadm, \tleq, \tjoin, \tatten, \tpromote, \tdemote)$ that assigns
calibrated trust to AI-generated evidence.  The key results are:

\begin{enumerate}[leftmargin=2em]
  \item The lattice satisfies the partial-order axioms and absorption laws
        (Theorems~\ref{thm:partial-order} and~\ref{thm:absorption}).
  \item Composition is associative (Theorem~\ref{thm:assoc}).
  \item The no-silent-promotion invariant prevents Copilot evidence from
        exceeding $\OracProp$ (Theorem~\ref{thm:no-silent}).
  \item Weighted composition with calibrated weights for each
        evidence source reflects empirical acceptance rates.
  \item Transport attenuation saturates at $\Unver$, preventing
        artificial contradiction.
\end{enumerate}

Across \expTotalPrograms{} programs, the algebra operates with sub-millisecond
mean composition time and zero axiom violations.

\paragraph{Future directions.}
We plan to extend the lattice to support \emph{probabilistic trust levels}
where $\CopSugg$ carries a continuous confidence score.  We also intend to
formalise the weighted composition in \LEAN~4 and investigate trust decay
models where Copilot evidence degrades over time as the codebase evolves.
A promising direction is \emph{contextual trust}, where the same evidence
source may receive different trust levels depending on the domain: Copilot
suggestions about string formatting may warrant higher trust than suggestions
about concurrent data structures.  We also plan to study the interaction
between trust attenuation and the Čech cohomology obstruction machinery
of Paper~3, since trust degradation along covering refinements mirrors
the cohomological descent condition.

\paragraph{Acknowledgments.}
We thank the Z3 team at Microsoft Research for their solver, the Lean
community for proof-assistant infrastructure, and the GitHub Copilot team
for LLM integration support.

\bibliographystyle{alpha}
\begin{thebibliography}{99}

\bibitem{jugeo-paper04}
\JuGeo{} Development Team.
\newblock Trust algebra for judgment-geometric verification.
\newblock Paper~4 in the \JuGeo{} series.

\bibitem{jugeo-paper01}
\JuGeo{} Development Team.
\newblock Semantic sites for program verification.
\newblock Paper~1 in the \JuGeo{} series.

\bibitem{jugeo-paper02}
\JuGeo{} Development Team.
\newblock Judgment algebra and the 8-tuple.
\newblock Paper~2 in the \JuGeo{} series.

\bibitem{jugeo-paper06}
\JuGeo{} Development Team.
\newblock Semantic moves in the judgment geometry.
\newblock Paper~6 in the \JuGeo{} series.

\bibitem{leino2010dafny}
K.~R.~M.~Leino.
\newblock Dafny: An automatic program verifier for functional correctness.
\newblock In \emph{LPAR}, 2010.

\bibitem{jung2018iris}
R.~Jung, R.~Krebbers, J.-H.~Jourdan, A.~Bizjak, L.~Birkedal, and
D.~Dreyer.
\newblock Iris from the ground up.
\newblock \emph{J.~Funct.~Program.}, 28:e20, 2018.

\bibitem{bell1973secure}
D.~E.~Bell and L.~J.~LaPadula.
\newblock Secure computer systems: Mathematical foundations.
\newblock Technical Report MTR-2547, MITRE, 1973.

\bibitem{li2003datalog}
N.~Li, B.~N.~Grosof, and J.~Feigenbaum.
\newblock Delegation logic: A logic-based approach to distributed authorization.
\newblock \emph{ACM Trans.~Inf.~Syst.~Secur.}, 6(1):128--171, 2003.

\bibitem{first2023baldur}
E.~First, M.~Rabe, T.~Ringer, and Y.~Brun.
\newblock Baldur: Whole-proof generation and repair with large language models.
\newblock In \emph{FSE}, 2023.

\bibitem{azerbayev2024llemma}
Z.~Azerbayev, H.~Schoelkopf, K.~Paster, et~al.
\newblock Llemma: An open language model for mathematics.
\newblock In \emph{ICLR}, 2024.

\bibitem{deepmind2024alphaproof}
DeepMind.
\newblock AlphaProof and AlphaGeometry~2.
\newblock Technical report, Google DeepMind, 2024.

\bibitem{harrison1996hol}
J.~Harrison.
\newblock {HOL Light}: A tutorial introduction.
\newblock In \emph{FMCAD}, 1996.

\bibitem{bertot2004interactive}
Y.~Bertot and P.~Cast{\'e}ran.
\newblock \emph{Interactive Theorem Proving and Program Development: Coq'Art}.
\newblock Springer, 2004.

\bibitem{goguen1992sheaves}
J.~A.~Goguen.
\newblock Sheaf semantics for concurrent interacting objects.
\newblock \emph{Math.~Struct.~Comput.~Sci.}, 2(2):159--191, 1992.

\bibitem{robinson2014topological}
M.~Robinson.
\newblock \emph{Topological Signal Processing}.
\newblock Springer, 2014.

\bibitem{curry2014sheaves}
J.~M.~Curry.
\newblock Sheaves, cosheaves and applications.
\newblock Ph.D. thesis, University of Pennsylvania, 2014.

\bibitem{hartshorne1977algebraic}
R.~Hartshorne.
\newblock \emph{Algebraic Geometry}.
\newblock Springer, 1977.

\bibitem{bott1982differential}
R.~Bott and L.~W.~Tu.
\newblock \emph{Differential Forms in Algebraic Topology}.
\newblock Springer, 1982.

\bibitem{demoura2008z3}
L.~de~Moura and N.~Bj{\o}rner.
\newblock {Z3}: An efficient {SMT} solver.
\newblock In \emph{TACAS}, 2008.

\bibitem{denning1976lattice}
D.~E.~Denning.
\newblock A lattice model of secure information flow.
\newblock \emph{Commun.~ACM}, 19(5):236--243, 1976.

\bibitem{mathur2021survey}
A.~P.~Mathur.
\newblock \emph{Foundations of Software Testing}.
\newblock Pearson, 2021.

\bibitem{yang2019survey}
Z.~Yang, C.~Qi, and Y.~Chen.
\newblock A survey of formal verification for smart contracts.
\newblock \emph{IEEE Access}, 7:164654--164679, 2019.

\bibitem{hales2017formal}
T.~C.~Hales et~al.
\newblock A formal proof of the {K}epler conjecture.
\newblock \emph{Forum Math. Pi}, 5:e2, 2017.

\bibitem{klein2009sel4}
G.~Klein et~al.
\newblock {seL4}: Formal verification of an {OS} kernel.
\newblock In \emph{SOSP}, 2009.

\bibitem{muller2016viper}
P.~M{\"u}ller, M.~Schwerhoff, and A.~J.~Summers.
\newblock Viper: A verification infrastructure for permission-based reasoning.
\newblock In \emph{VMCAI}, 2016.

\bibitem{chen2021evaluating}
M.~Chen, J.~Tworek, H.~Jun, et~al.
\newblock Evaluating large language models trained on code.
\newblock \emph{arXiv:2107.03374}, 2021.

\bibitem{austin2021program}
J.~Austin, A.~Odena, M.~Nye, et~al.
\newblock Program synthesis with large language models.
\newblock \emph{arXiv:2108.07732}, 2021.

\bibitem{ouyang2022training}
L.~Ouyang, J.~Wu, X.~Jiang, et~al.
\newblock Training language models to follow instructions with human feedback.
\newblock In \emph{NeurIPS}, 2022.

\bibitem{davila2023lattice}
J.~D{\'a}vila.
\newblock Lattice-based trust for distributed systems.
\newblock In \emph{FORTE}, 2023.

\bibitem{blazy2006formal}
S.~Blazy, Z.~Dargaye, and X.~Leroy.
\newblock Formal verification of a {C} compiler front-end.
\newblock In \emph{FM}, 2006.

\bibitem{nipkow2002isabelle}
T.~Nipkow, L.~C.~Paulson, and M.~Wenzel.
\newblock \emph{Isabelle/HOL: A Proof Assistant for Higher-Order Logic}.
\newblock Springer, 2002.

\bibitem{mathlib2020}
The mathlib Community.
\newblock The {L}ean mathematical library.
\newblock In \emph{CPP}, 2020.

\bibitem{swamy2016dependent}
N.~Swamy, C.~Hri\c{t}cu, C.~Keller, et~al.
\newblock Dependent types and multi-monadic effects in {F*}.
\newblock In \emph{POPL}, 2016.

\end{thebibliography}

\end{document}
